\documentclass[11pt]{article}
\usepackage[margin=0.85in]{geometry}
\usepackage[T1]{fontenc}
\usepackage{lmodern}
\usepackage{amsmath,amssymb,amsthm,mathtools,booktabs,longtable,array}
\usepackage{listings,xcolor,xurl,hyperref}
\setlength{\emergencystretch}{2em}
\hypersetup{colorlinks=true,linkcolor=blue,urlcolor=blue,citecolor=blue,
 pdftitle={Distance loss in girth-eight cyclic lifted-product codes},pdfauthor={}}
\newtheorem{theorem}{Theorem}[section]
\newtheorem{lemma}[theorem]{Lemma}
\newtheorem{proposition}[theorem]{Proposition}
\newtheorem{corollary}[theorem]{Corollary}
\theoremstyle{definition}\newtheorem{definition}[theorem]{Definition}
\newcommand{\F}{\mathbb F_2}
\newcommand{\Z}{\mathbb Z}
\newcommand{\wt}{\operatorname{wt}}
\newcommand{\row}{\operatorname{row}}
\newcommand{\HX}{H_{\mathsf X}}
\newcommand{\HZ}{H_{\mathsf Z}}
\newcommand{\proofcommit}{98eb91cd9049e6a88887a5ba79150ad96a4a3e53}
\newcommand{\proofbase}{https://github.com/alectionary/vibe-math/blob/\proofcommit/formalization/distance-preservation}
% Generated declaration line anchors at the immutable proof commit.
\expandafter\def\csname leanline@BinaryFinite.countP_ofFn\endcsname{18}
\expandafter\def\csname leanline@BinaryFinite.cpop_list\endcsname{10}
\expandafter\def\csname leanline@BinaryFinite.pack44\endcsname{7}
\expandafter\def\csname leanline@BinaryFinite.pack44_bit\endcsname{8}
\expandafter\def\csname leanline@BinaryFinite.pack44_weight\endcsname{26}
\expandafter\def\csname leanline@BinaryFinite.xor_false_iff_even\endcsname{31}
\expandafter\def\csname leanline@BinaryWords.InRowSpan\endcsname{72}
\expandafter\def\csname leanline@BinaryWords.dot\endcsname{20}
\expandafter\def\csname leanline@BinaryWords.dotOn\endcsname{17}
\expandafter\def\csname leanline@BinaryWords.dotOn_comm\endcsname{54}
\expandafter\def\csname leanline@BinaryWords.dotOn_xor\endcsname{26}
\expandafter\def\csname leanline@BinaryWords.dotOn_zero\endcsname{43}
\expandafter\def\csname leanline@BinaryWords.dot_comm\endcsname{63}
\expandafter\def\csname leanline@BinaryWords.dot_rowSum\endcsname{75}
\expandafter\def\csname leanline@BinaryWords.dot_xor\endcsname{40}
\expandafter\def\csname leanline@BinaryWords.dot_xor_left\endcsname{65}
\expandafter\def\csname leanline@BinaryWords.dot_zero\endcsname{52}
\expandafter\def\csname leanline@BinaryWords.equation_of_row_relation\endcsname{93}
\expandafter\def\csname leanline@BinaryWords.forall_fin_of_all\endcsname{6}
\expandafter\def\csname leanline@BinaryWords.forall_fin_two_of_all\endcsname{11}
\expandafter\def\csname leanline@BinaryWords.getLsbD_bool_ite\endcsname{22}
\expandafter\def\csname leanline@BinaryWords.not_in_row_span\endcsname{85}
\expandafter\def\csname leanline@BinaryWords.rowSum\endcsname{69}
\expandafter\def\csname leanline@ClassicalCertificate11.actual_classical_lower\endcsname{88}
\expandafter\def\csname leanline@ClassicalCertificate11.all_blocks\endcsname{8}
\expandafter\def\csname leanline@ClassicalCertificate11.block_0\endcsname{22}
\expandafter\def\csname leanline@ClassicalCertificate11.block_1\endcsname{23}
\expandafter\def\csname leanline@ClassicalCertificate11.block_10\endcsname{32}
\expandafter\def\csname leanline@ClassicalCertificate11.block_11\endcsname{33}
\expandafter\def\csname leanline@ClassicalCertificate11.block_12\endcsname{34}
\expandafter\def\csname leanline@ClassicalCertificate11.block_13\endcsname{35}
\expandafter\def\csname leanline@ClassicalCertificate11.block_14\endcsname{36}
\expandafter\def\csname leanline@ClassicalCertificate11.block_15\endcsname{37}
\expandafter\def\csname leanline@ClassicalCertificate11.block_16\endcsname{38}
\expandafter\def\csname leanline@ClassicalCertificate11.block_17\endcsname{39}
\expandafter\def\csname leanline@ClassicalCertificate11.block_18\endcsname{40}
\expandafter\def\csname leanline@ClassicalCertificate11.block_19\endcsname{41}
\expandafter\def\csname leanline@ClassicalCertificate11.block_2\endcsname{24}
\expandafter\def\csname leanline@ClassicalCertificate11.block_20\endcsname{42}
\expandafter\def\csname leanline@ClassicalCertificate11.block_21\endcsname{43}
\expandafter\def\csname leanline@ClassicalCertificate11.block_22\endcsname{44}
\expandafter\def\csname leanline@ClassicalCertificate11.block_23\endcsname{45}
\expandafter\def\csname leanline@ClassicalCertificate11.block_24\endcsname{46}
\expandafter\def\csname leanline@ClassicalCertificate11.block_25\endcsname{47}
\expandafter\def\csname leanline@ClassicalCertificate11.block_26\endcsname{48}
\expandafter\def\csname leanline@ClassicalCertificate11.block_27\endcsname{49}
\expandafter\def\csname leanline@ClassicalCertificate11.block_28\endcsname{50}
\expandafter\def\csname leanline@ClassicalCertificate11.block_29\endcsname{51}
\expandafter\def\csname leanline@ClassicalCertificate11.block_3\endcsname{25}
\expandafter\def\csname leanline@ClassicalCertificate11.block_30\endcsname{52}
\expandafter\def\csname leanline@ClassicalCertificate11.block_31\endcsname{53}
\expandafter\def\csname leanline@ClassicalCertificate11.block_32\endcsname{54}
\expandafter\def\csname leanline@ClassicalCertificate11.block_33\endcsname{55}
\expandafter\def\csname leanline@ClassicalCertificate11.block_34\endcsname{56}
\expandafter\def\csname leanline@ClassicalCertificate11.block_35\endcsname{57}
\expandafter\def\csname leanline@ClassicalCertificate11.block_36\endcsname{58}
\expandafter\def\csname leanline@ClassicalCertificate11.block_37\endcsname{59}
\expandafter\def\csname leanline@ClassicalCertificate11.block_38\endcsname{60}
\expandafter\def\csname leanline@ClassicalCertificate11.block_39\endcsname{61}
\expandafter\def\csname leanline@ClassicalCertificate11.block_4\endcsname{26}
\expandafter\def\csname leanline@ClassicalCertificate11.block_40\endcsname{62}
\expandafter\def\csname leanline@ClassicalCertificate11.block_41\endcsname{63}
\expandafter\def\csname leanline@ClassicalCertificate11.block_42\endcsname{64}
\expandafter\def\csname leanline@ClassicalCertificate11.block_43\endcsname{65}
\expandafter\def\csname leanline@ClassicalCertificate11.block_44\endcsname{66}
\expandafter\def\csname leanline@ClassicalCertificate11.block_45\endcsname{67}
\expandafter\def\csname leanline@ClassicalCertificate11.block_46\endcsname{68}
\expandafter\def\csname leanline@ClassicalCertificate11.block_47\endcsname{69}
\expandafter\def\csname leanline@ClassicalCertificate11.block_48\endcsname{70}
\expandafter\def\csname leanline@ClassicalCertificate11.block_49\endcsname{71}
\expandafter\def\csname leanline@ClassicalCertificate11.block_5\endcsname{27}
\expandafter\def\csname leanline@ClassicalCertificate11.block_50\endcsname{72}
\expandafter\def\csname leanline@ClassicalCertificate11.block_51\endcsname{73}
\expandafter\def\csname leanline@ClassicalCertificate11.block_52\endcsname{74}
\expandafter\def\csname leanline@ClassicalCertificate11.block_53\endcsname{75}
\expandafter\def\csname leanline@ClassicalCertificate11.block_54\endcsname{76}
\expandafter\def\csname leanline@ClassicalCertificate11.block_55\endcsname{77}
\expandafter\def\csname leanline@ClassicalCertificate11.block_56\endcsname{78}
\expandafter\def\csname leanline@ClassicalCertificate11.block_57\endcsname{79}
\expandafter\def\csname leanline@ClassicalCertificate11.block_58\endcsname{80}
\expandafter\def\csname leanline@ClassicalCertificate11.block_59\endcsname{81}
\expandafter\def\csname leanline@ClassicalCertificate11.block_6\endcsname{28}
\expandafter\def\csname leanline@ClassicalCertificate11.block_60\endcsname{82}
\expandafter\def\csname leanline@ClassicalCertificate11.block_61\endcsname{83}
\expandafter\def\csname leanline@ClassicalCertificate11.block_62\endcsname{84}
\expandafter\def\csname leanline@ClassicalCertificate11.block_63\endcsname{85}
\expandafter\def\csname leanline@ClassicalCertificate11.block_7\endcsname{29}
\expandafter\def\csname leanline@ClassicalCertificate11.block_8\endcsname{30}
\expandafter\def\csname leanline@ClassicalCertificate11.block_9\endcsname{31}
\expandafter\def\csname leanline@ClassicalCertificate11.encode\endcsname{5}
\expandafter\def\csname leanline@ClassicalCertificate11.encoded_lower\endcsname{75}
\expandafter\def\csname leanline@ClassicalCertificate11.goodBlock\endcsname{19}
\expandafter\def\csname leanline@ClassicalUpper20.classical_distance_upper\endcsname{62}
\expandafter\def\csname leanline@ClassicalUpper20.kernel\endcsname{14}
\expandafter\def\csname leanline@ClassicalUpper20.nonzero\endcsname{20}
\expandafter\def\csname leanline@ClassicalUpper20.supportList\endcsname{41}
\expandafter\def\csname leanline@ClassicalUpper20.weight_le\endcsname{50}
\expandafter\def\csname leanline@ClassicalUpper20.word\endcsname{8}
\expandafter\def\csname leanline@ClassicalUpper20.zero_outside\endcsname{46}
\expandafter\def\csname leanline@CoverDistance.lower_bound\endcsname{9}
\expandafter\def\csname leanline@CoverWeight.nonzero_one\endcsname{13}
\expandafter\def\csname leanline@CoverWeight.push\endcsname{10}
\expandafter\def\csname leanline@CoverWeight.push_eq_card\endcsname{28}
\expandafter\def\csname leanline@CoverWeight.push_support_subset\endcsname{37}
\expandafter\def\csname leanline@CoverWeight.push_weight_le\endcsname{51}
\expandafter\def\csname leanline@CoverWeight.sum_eq_card_support\endcsname{18}
\expandafter\def\csname leanline@CoverWeight.support\endcsname{9}
\expandafter\def\csname leanline@CoverWeight.zero_push_fiber_lower\endcsname{55}
\expandafter\def\csname leanline@CoverWeight.zero_push_weight_doubles\endcsname{72}
\expandafter\def\csname leanline@CyclicCover.Adjacent\endcsname{15}
\expandafter\def\csname leanline@CyclicCover.Kernel\endcsname{14}
\expandafter\def\csname leanline@CyclicCover.Word\endcsname{10}
\expandafter\def\csname leanline@CyclicCover.boundary\endcsname{12}
\expandafter\def\csname leanline@CyclicCover.boundary_push\endcsname{34}
\expandafter\def\csname leanline@CyclicCover.fiber_shift\endcsname{26}
\expandafter\def\csname leanline@CyclicCover.kernel_push\endcsname{44}
\expandafter\def\csname leanline@CyclicCover.project\endcsname{11}
\expandafter\def\csname leanline@CyclicCover.projected_support_stopping\endcsname{67}
\expandafter\def\csname leanline@CyclicCover.push_apply\endcsname{20}
\expandafter\def\csname leanline@CyclicCover.stopping_of_partners\endcsname{51}
\expandafter\def\csname leanline@DistanceFamily.ClassicalHasWeight\endcsname{64}
\expandafter\def\csname leanline@DistanceFamily.LogicalX\endcsname{39}
\expandafter\def\csname leanline@DistanceFamily.QuantumHasWeight\endcsname{69}
\expandafter\def\csname leanline@DistanceFamily.classicalDistance\endcsname{81}
\expandafter\def\csname leanline@DistanceFamily.classical_distance_lower\endcsname{84}
\expandafter\def\csname leanline@DistanceFamily.classical_distance_upper\endcsname{16}
\expandafter\def\csname leanline@DistanceFamily.classical_exists\endcsname{77}
\expandafter\def\csname leanline@DistanceFamily.classical_lower_twelve\endcsname{11}
\expandafter\def\csname leanline@DistanceFamily.distance_drop\endcsname{97}
\expandafter\def\csname leanline@DistanceFamily.explicit_logical\endcsname{44}
\expandafter\def\csname leanline@DistanceFamily.explicit_weight\endcsname{49}
\expandafter\def\csname leanline@DistanceFamily.fiberPair\endcsname{53}
\expandafter\def\csname leanline@DistanceFamily.fiberPair_kernel\endcsname{54}
\expandafter\def\csname leanline@DistanceFamily.fiberPair_nonzero\endcsname{59}
\expandafter\def\csname leanline@DistanceFamily.fiberPair_weight\endcsname{7}
\expandafter\def\csname leanline@DistanceFamily.infinite_counterfamily\endcsname{102}
\expandafter\def\csname leanline@DistanceFamily.join\endcsname{35}
\expandafter\def\csname leanline@DistanceFamily.liftE\endcsname{9}
\expandafter\def\csname leanline@DistanceFamily.quantumDistance\endcsname{82}
\expandafter\def\csname leanline@DistanceFamily.quantum_distance_upper\endcsname{89}
\expandafter\def\csname leanline@DistanceFamily.quantum_exists\endcsname{79}
\expandafter\def\csname leanline@DistanceFamily.stronger_counterfamily\endcsname{28}
\expandafter\def\csname leanline@DistanceFamily.transpose\endcsname{38}
\expandafter\def\csname leanline@DistanceFamily.transpose_c_action\endcsname{22}
\expandafter\def\csname leanline@FamilyBase11.ClassicalKernel\endcsname{12}
\expandafter\def\csname leanline@FamilyBase11.bits\endcsname{10}
\expandafter\def\csname leanline@FamilyBase11.classicalRow\endcsname{11}
\expandafter\def\csname leanline@FamilyBase11.dot_row\endcsname{106}
\expandafter\def\csname leanline@FamilyBase11.dot_row_0\endcsname{7}
\expandafter\def\csname leanline@FamilyBase11.dot_row_1\endcsname{10}
\expandafter\def\csname leanline@FamilyBase11.dot_row_10\endcsname{37}
\expandafter\def\csname leanline@FamilyBase11.dot_row_11\endcsname{40}
\expandafter\def\csname leanline@FamilyBase11.dot_row_12\endcsname{43}
\expandafter\def\csname leanline@FamilyBase11.dot_row_13\endcsname{46}
\expandafter\def\csname leanline@FamilyBase11.dot_row_14\endcsname{49}
\expandafter\def\csname leanline@FamilyBase11.dot_row_15\endcsname{52}
\expandafter\def\csname leanline@FamilyBase11.dot_row_16\endcsname{55}
\expandafter\def\csname leanline@FamilyBase11.dot_row_17\endcsname{58}
\expandafter\def\csname leanline@FamilyBase11.dot_row_18\endcsname{61}
\expandafter\def\csname leanline@FamilyBase11.dot_row_19\endcsname{64}
\expandafter\def\csname leanline@FamilyBase11.dot_row_2\endcsname{13}
\expandafter\def\csname leanline@FamilyBase11.dot_row_20\endcsname{67}
\expandafter\def\csname leanline@FamilyBase11.dot_row_21\endcsname{70}
\expandafter\def\csname leanline@FamilyBase11.dot_row_22\endcsname{73}
\expandafter\def\csname leanline@FamilyBase11.dot_row_23\endcsname{76}
\expandafter\def\csname leanline@FamilyBase11.dot_row_24\endcsname{79}
\expandafter\def\csname leanline@FamilyBase11.dot_row_25\endcsname{82}
\expandafter\def\csname leanline@FamilyBase11.dot_row_26\endcsname{85}
\expandafter\def\csname leanline@FamilyBase11.dot_row_27\endcsname{88}
\expandafter\def\csname leanline@FamilyBase11.dot_row_28\endcsname{91}
\expandafter\def\csname leanline@FamilyBase11.dot_row_29\endcsname{94}
\expandafter\def\csname leanline@FamilyBase11.dot_row_3\endcsname{16}
\expandafter\def\csname leanline@FamilyBase11.dot_row_30\endcsname{97}
\expandafter\def\csname leanline@FamilyBase11.dot_row_31\endcsname{100}
\expandafter\def\csname leanline@FamilyBase11.dot_row_32\endcsname{103}
\expandafter\def\csname leanline@FamilyBase11.dot_row_4\endcsname{19}
\expandafter\def\csname leanline@FamilyBase11.dot_row_5\endcsname{22}
\expandafter\def\csname leanline@FamilyBase11.dot_row_6\endcsname{25}
\expandafter\def\csname leanline@FamilyBase11.dot_row_7\endcsname{28}
\expandafter\def\csname leanline@FamilyBase11.dot_row_8\endcsname{31}
\expandafter\def\csname leanline@FamilyBase11.dot_row_9\endcsname{34}
\expandafter\def\csname leanline@FamilyBase11.encode_bit_0\endcsname{238}
\expandafter\def\csname leanline@FamilyBase11.encode_bit_1\endcsname{242}
\expandafter\def\csname leanline@FamilyBase11.encode_bit_10\endcsname{278}
\expandafter\def\csname leanline@FamilyBase11.encode_bit_11\endcsname{282}
\expandafter\def\csname leanline@FamilyBase11.encode_bit_12\endcsname{286}
\expandafter\def\csname leanline@FamilyBase11.encode_bit_13\endcsname{290}
\expandafter\def\csname leanline@FamilyBase11.encode_bit_14\endcsname{294}
\expandafter\def\csname leanline@FamilyBase11.encode_bit_15\endcsname{298}
\expandafter\def\csname leanline@FamilyBase11.encode_bit_16\endcsname{302}
\expandafter\def\csname leanline@FamilyBase11.encode_bit_17\endcsname{306}
\expandafter\def\csname leanline@FamilyBase11.encode_bit_18\endcsname{310}
\expandafter\def\csname leanline@FamilyBase11.encode_bit_19\endcsname{314}
\expandafter\def\csname leanline@FamilyBase11.encode_bit_2\endcsname{246}
\expandafter\def\csname leanline@FamilyBase11.encode_bit_20\endcsname{318}
\expandafter\def\csname leanline@FamilyBase11.encode_bit_21\endcsname{322}
\expandafter\def\csname leanline@FamilyBase11.encode_bit_22\endcsname{326}
\expandafter\def\csname leanline@FamilyBase11.encode_bit_23\endcsname{330}
\expandafter\def\csname leanline@FamilyBase11.encode_bit_24\endcsname{334}
\expandafter\def\csname leanline@FamilyBase11.encode_bit_25\endcsname{338}
\expandafter\def\csname leanline@FamilyBase11.encode_bit_26\endcsname{342}
\expandafter\def\csname leanline@FamilyBase11.encode_bit_27\endcsname{346}
\expandafter\def\csname leanline@FamilyBase11.encode_bit_28\endcsname{350}
\expandafter\def\csname leanline@FamilyBase11.encode_bit_29\endcsname{354}
\expandafter\def\csname leanline@FamilyBase11.encode_bit_3\endcsname{250}
\expandafter\def\csname leanline@FamilyBase11.encode_bit_30\endcsname{358}
\expandafter\def\csname leanline@FamilyBase11.encode_bit_31\endcsname{362}
\expandafter\def\csname leanline@FamilyBase11.encode_bit_32\endcsname{366}
\expandafter\def\csname leanline@FamilyBase11.encode_bit_33\endcsname{370}
\expandafter\def\csname leanline@FamilyBase11.encode_bit_34\endcsname{374}
\expandafter\def\csname leanline@FamilyBase11.encode_bit_35\endcsname{378}
\expandafter\def\csname leanline@FamilyBase11.encode_bit_36\endcsname{382}
\expandafter\def\csname leanline@FamilyBase11.encode_bit_37\endcsname{386}
\expandafter\def\csname leanline@FamilyBase11.encode_bit_38\endcsname{390}
\expandafter\def\csname leanline@FamilyBase11.encode_bit_39\endcsname{394}
\expandafter\def\csname leanline@FamilyBase11.encode_bit_4\endcsname{254}
\expandafter\def\csname leanline@FamilyBase11.encode_bit_40\endcsname{398}
\expandafter\def\csname leanline@FamilyBase11.encode_bit_41\endcsname{402}
\expandafter\def\csname leanline@FamilyBase11.encode_bit_42\endcsname{406}
\expandafter\def\csname leanline@FamilyBase11.encode_bit_43\endcsname{410}
\expandafter\def\csname leanline@FamilyBase11.encode_bit_5\endcsname{258}
\expandafter\def\csname leanline@FamilyBase11.encode_bit_6\endcsname{262}
\expandafter\def\csname leanline@FamilyBase11.encode_bit_7\endcsname{266}
\expandafter\def\csname leanline@FamilyBase11.encode_bit_8\endcsname{270}
\expandafter\def\csname leanline@FamilyBase11.encode_bit_9\endcsname{274}
\expandafter\def\csname leanline@FamilyBase11.exponent\endcsname{9}
\expandafter\def\csname leanline@FamilyBase11.exponents\endcsname{8}
\expandafter\def\csname leanline@FamilyBase11.freeInput\endcsname{414}
\expandafter\def\csname leanline@FamilyBase11.kernel_reconstruct\endcsname{415}
\expandafter\def\csname leanline@FamilyBase11.mask_1\endcsname{13}
\expandafter\def\csname leanline@FamilyBase11.mask_10002978832384\endcsname{190}
\expandafter\def\csname leanline@FamilyBase11.mask_1024\endcsname{43}
\expandafter\def\csname leanline@FamilyBase11.mask_1048576\endcsname{73}
\expandafter\def\csname leanline@FamilyBase11.mask_1073741824\endcsname{103}
\expandafter\def\csname leanline@FamilyBase11.mask_10741715304448\endcsname{193}
\expandafter\def\csname leanline@FamilyBase11.mask_1099511627776\endcsname{139}
\expandafter\def\csname leanline@FamilyBase11.mask_11712375816192\endcsname{196}
\expandafter\def\csname leanline@FamilyBase11.mask_12644385816576\endcsname{199}
\expandafter\def\csname leanline@FamilyBase11.mask_128\endcsname{34}
\expandafter\def\csname leanline@FamilyBase11.mask_13039522807808\endcsname{202}
\expandafter\def\csname leanline@FamilyBase11.mask_131072\endcsname{64}
\expandafter\def\csname leanline@FamilyBase11.mask_134217728\endcsname{94}
\expandafter\def\csname leanline@FamilyBase11.mask_13632228294656\endcsname{205}
\expandafter\def\csname leanline@FamilyBase11.mask_137438953472\endcsname{124}
\expandafter\def\csname leanline@FamilyBase11.mask_14117559599104\endcsname{208}
\expandafter\def\csname leanline@FamilyBase11.mask_14147624370176\endcsname{211}
\expandafter\def\csname leanline@FamilyBase11.mask_14220636717056\endcsname{214}
\expandafter\def\csname leanline@FamilyBase11.mask_14830522073088\endcsname{217}
\expandafter\def\csname leanline@FamilyBase11.mask_15337328214016\endcsname{220}
\expandafter\def\csname leanline@FamilyBase11.mask_15358805147648\endcsname{223}
\expandafter\def\csname leanline@FamilyBase11.mask_15380279984128\endcsname{226}
\expandafter\def\csname leanline@FamilyBase11.mask_16\endcsname{25}
\expandafter\def\csname leanline@FamilyBase11.mask_16316580757504\endcsname{229}
\expandafter\def\csname leanline@FamilyBase11.mask_16384\endcsname{55}
\expandafter\def\csname leanline@FamilyBase11.mask_16595755728896\endcsname{232}
\expandafter\def\csname leanline@FamilyBase11.mask_16777216\endcsname{85}
\expandafter\def\csname leanline@FamilyBase11.mask_16900698406912\endcsname{235}
\expandafter\def\csname leanline@FamilyBase11.mask_17179869184\endcsname{115}
\expandafter\def\csname leanline@FamilyBase11.mask_197570592768\endcsname{127}
\expandafter\def\csname leanline@FamilyBase11.mask_2\endcsname{16}
\expandafter\def\csname leanline@FamilyBase11.mask_2048\endcsname{46}
\expandafter\def\csname leanline@FamilyBase11.mask_2097152\endcsname{76}
\expandafter\def\csname leanline@FamilyBase11.mask_2147483648\endcsname{106}
\expandafter\def\csname leanline@FamilyBase11.mask_2199023255552\endcsname{142}
\expandafter\def\csname leanline@FamilyBase11.mask_2503965933568\endcsname{145}
\expandafter\def\csname leanline@FamilyBase11.mask_256\endcsname{37}
\expandafter\def\csname leanline@FamilyBase11.mask_2619932147712\endcsname{148}
\expandafter\def\csname leanline@FamilyBase11.mask_262144\endcsname{67}
\expandafter\def\csname leanline@FamilyBase11.mask_2658586853376\endcsname{151}
\expandafter\def\csname leanline@FamilyBase11.mask_268435456\endcsname{97}
\expandafter\def\csname leanline@FamilyBase11.mask_274877906944\endcsname{130}
\expandafter\def\csname leanline@FamilyBase11.mask_2787435872256\endcsname{154}
\expandafter\def\csname leanline@FamilyBase11.mask_32\endcsname{28}
\expandafter\def\csname leanline@FamilyBase11.mask_32768\endcsname{58}
\expandafter\def\csname leanline@FamilyBase11.mask_33554432\endcsname{88}
\expandafter\def\csname leanline@FamilyBase11.mask_3375846391808\endcsname{157}
\expandafter\def\csname leanline@FamilyBase11.mask_34359738368\endcsname{118}
\expandafter\def\csname leanline@FamilyBase11.mask_3869765533696\endcsname{160}
\expandafter\def\csname leanline@FamilyBase11.mask_4\endcsname{19}
\expandafter\def\csname leanline@FamilyBase11.mask_4096\endcsname{49}
\expandafter\def\csname leanline@FamilyBase11.mask_4194304\endcsname{79}
\expandafter\def\csname leanline@FamilyBase11.mask_4294967296\endcsname{109}
\expandafter\def\csname leanline@FamilyBase11.mask_4398046511104\endcsname{163}
\expandafter\def\csname leanline@FamilyBase11.mask_512\endcsname{40}
\expandafter\def\csname leanline@FamilyBase11.mask_524288\endcsname{70}
\expandafter\def\csname leanline@FamilyBase11.mask_536870912\endcsname{100}
\expandafter\def\csname leanline@FamilyBase11.mask_549755813888\endcsname{133}
\expandafter\def\csname leanline@FamilyBase11.mask_5750961209344\endcsname{166}
\expandafter\def\csname leanline@FamilyBase11.mask_6025841213440\endcsname{169}
\expandafter\def\csname leanline@FamilyBase11.mask_639952224256\endcsname{136}
\expandafter\def\csname leanline@FamilyBase11.mask_64\endcsname{31}
\expandafter\def\csname leanline@FamilyBase11.mask_6519762452480\endcsname{172}
\expandafter\def\csname leanline@FamilyBase11.mask_65536\endcsname{61}
\expandafter\def\csname leanline@FamilyBase11.mask_67108864\endcsname{91}
\expandafter\def\csname leanline@FamilyBase11.mask_6738805784576\endcsname{175}
\expandafter\def\csname leanline@FamilyBase11.mask_68719476736\endcsname{121}
\expandafter\def\csname leanline@FamilyBase11.mask_7735236100096\endcsname{178}
\expandafter\def\csname leanline@FamilyBase11.mask_7898446954496\endcsname{181}
\expandafter\def\csname leanline@FamilyBase11.mask_8\endcsname{22}
\expandafter\def\csname leanline@FamilyBase11.mask_8192\endcsname{52}
\expandafter\def\csname leanline@FamilyBase11.mask_8388608\endcsname{82}
\expandafter\def\csname leanline@FamilyBase11.mask_8589934592\endcsname{112}
\expandafter\def\csname leanline@FamilyBase11.mask_8796093022208\endcsname{184}
\expandafter\def\csname leanline@FamilyBase11.mask_8800390086656\endcsname{187}
\expandafter\def\csname leanline@FamilyBaseGraph.Check\endcsname{9}
\expandafter\def\csname leanline@FamilyBaseGraph.E\endcsname{13}
\expandafter\def\csname leanline@FamilyBaseGraph.Var\endcsname{8}
\expandafter\def\csname leanline@FamilyBaseGraph.bool_four\endcsname{14}
\expandafter\def\csname leanline@FamilyBaseGraph.boundary\endcsname{16}
\expandafter\def\csname leanline@FamilyBaseGraph.checkEquiv\endcsname{11}
\expandafter\def\csname leanline@FamilyBaseGraph.code_lower_twelve\endcsname{321}
\expandafter\def\csname leanline@FamilyBaseGraph.e\endcsname{12}
\expandafter\def\csname leanline@FamilyBaseGraph.incidence_bridge\endcsname{16}
\expandafter\def\csname leanline@FamilyBaseGraph.kernel_row\endcsname{281}
\expandafter\def\csname leanline@FamilyBaseGraph.kernel_row_0\endcsname{17}
\expandafter\def\csname leanline@FamilyBaseGraph.kernel_row_1\endcsname{25}
\expandafter\def\csname leanline@FamilyBaseGraph.kernel_row_10\endcsname{97}
\expandafter\def\csname leanline@FamilyBaseGraph.kernel_row_11\endcsname{105}
\expandafter\def\csname leanline@FamilyBaseGraph.kernel_row_12\endcsname{113}
\expandafter\def\csname leanline@FamilyBaseGraph.kernel_row_13\endcsname{121}
\expandafter\def\csname leanline@FamilyBaseGraph.kernel_row_14\endcsname{129}
\expandafter\def\csname leanline@FamilyBaseGraph.kernel_row_15\endcsname{137}
\expandafter\def\csname leanline@FamilyBaseGraph.kernel_row_16\endcsname{145}
\expandafter\def\csname leanline@FamilyBaseGraph.kernel_row_17\endcsname{153}
\expandafter\def\csname leanline@FamilyBaseGraph.kernel_row_18\endcsname{161}
\expandafter\def\csname leanline@FamilyBaseGraph.kernel_row_19\endcsname{169}
\expandafter\def\csname leanline@FamilyBaseGraph.kernel_row_2\endcsname{33}
\expandafter\def\csname leanline@FamilyBaseGraph.kernel_row_20\endcsname{177}
\expandafter\def\csname leanline@FamilyBaseGraph.kernel_row_21\endcsname{185}
\expandafter\def\csname leanline@FamilyBaseGraph.kernel_row_22\endcsname{193}
\expandafter\def\csname leanline@FamilyBaseGraph.kernel_row_23\endcsname{201}
\expandafter\def\csname leanline@FamilyBaseGraph.kernel_row_24\endcsname{209}
\expandafter\def\csname leanline@FamilyBaseGraph.kernel_row_25\endcsname{217}
\expandafter\def\csname leanline@FamilyBaseGraph.kernel_row_26\endcsname{225}
\expandafter\def\csname leanline@FamilyBaseGraph.kernel_row_27\endcsname{233}
\expandafter\def\csname leanline@FamilyBaseGraph.kernel_row_28\endcsname{241}
\expandafter\def\csname leanline@FamilyBaseGraph.kernel_row_29\endcsname{249}
\expandafter\def\csname leanline@FamilyBaseGraph.kernel_row_3\endcsname{41}
\expandafter\def\csname leanline@FamilyBaseGraph.kernel_row_30\endcsname{257}
\expandafter\def\csname leanline@FamilyBaseGraph.kernel_row_31\endcsname{265}
\expandafter\def\csname leanline@FamilyBaseGraph.kernel_row_32\endcsname{273}
\expandafter\def\csname leanline@FamilyBaseGraph.kernel_row_4\endcsname{49}
\expandafter\def\csname leanline@FamilyBaseGraph.kernel_row_5\endcsname{57}
\expandafter\def\csname leanline@FamilyBaseGraph.kernel_row_6\endcsname{65}
\expandafter\def\csname leanline@FamilyBaseGraph.kernel_row_7\endcsname{73}
\expandafter\def\csname leanline@FamilyBaseGraph.kernel_row_8\endcsname{81}
\expandafter\def\csname leanline@FamilyBaseGraph.kernel_row_9\endcsname{89}
\expandafter\def\csname leanline@FamilyBaseGraph.packWord\endcsname{9}
\expandafter\def\csname leanline@FamilyBaseGraph.packWord_bit\endcsname{10}
\expandafter\def\csname leanline@FamilyBaseGraph.packWord_bit_nat\endcsname{11}
\expandafter\def\csname leanline@FamilyBaseGraph.packWord_weight\endcsname{316}
\expandafter\def\csname leanline@FamilyBaseGraph.stopping_lower_six\endcsname{42}
\expandafter\def\csname leanline@FamilyBaseGraph.stopping_transport\endcsname{21}
\expandafter\def\csname leanline@FamilyBaseGraph.varEquiv\endcsname{10}
\expandafter\def\csname leanline@FamilyClassical11.actual_classical_lower\endcsname{88}
\expandafter\def\csname leanline@FamilyClassical11.all_blocks\endcsname{8}
\expandafter\def\csname leanline@FamilyClassical11.block_0\endcsname{7}
\expandafter\def\csname leanline@FamilyClassical11.block_1\endcsname{8}
\expandafter\def\csname leanline@FamilyClassical11.block_10\endcsname{17}
\expandafter\def\csname leanline@FamilyClassical11.block_11\endcsname{18}
\expandafter\def\csname leanline@FamilyClassical11.block_12\endcsname{19}
\expandafter\def\csname leanline@FamilyClassical11.block_13\endcsname{20}
\expandafter\def\csname leanline@FamilyClassical11.block_14\endcsname{21}
\expandafter\def\csname leanline@FamilyClassical11.block_15\endcsname{22}
\expandafter\def\csname leanline@FamilyClassical11.block_16\endcsname{23}
\expandafter\def\csname leanline@FamilyClassical11.block_17\endcsname{24}
\expandafter\def\csname leanline@FamilyClassical11.block_18\endcsname{25}
\expandafter\def\csname leanline@FamilyClassical11.block_19\endcsname{26}
\expandafter\def\csname leanline@FamilyClassical11.block_2\endcsname{9}
\expandafter\def\csname leanline@FamilyClassical11.block_20\endcsname{27}
\expandafter\def\csname leanline@FamilyClassical11.block_21\endcsname{28}
\expandafter\def\csname leanline@FamilyClassical11.block_22\endcsname{29}
\expandafter\def\csname leanline@FamilyClassical11.block_23\endcsname{30}
\expandafter\def\csname leanline@FamilyClassical11.block_24\endcsname{31}
\expandafter\def\csname leanline@FamilyClassical11.block_25\endcsname{32}
\expandafter\def\csname leanline@FamilyClassical11.block_26\endcsname{33}
\expandafter\def\csname leanline@FamilyClassical11.block_27\endcsname{34}
\expandafter\def\csname leanline@FamilyClassical11.block_28\endcsname{35}
\expandafter\def\csname leanline@FamilyClassical11.block_29\endcsname{36}
\expandafter\def\csname leanline@FamilyClassical11.block_3\endcsname{10}
\expandafter\def\csname leanline@FamilyClassical11.block_30\endcsname{37}
\expandafter\def\csname leanline@FamilyClassical11.block_31\endcsname{38}
\expandafter\def\csname leanline@FamilyClassical11.block_32\endcsname{39}
\expandafter\def\csname leanline@FamilyClassical11.block_33\endcsname{40}
\expandafter\def\csname leanline@FamilyClassical11.block_34\endcsname{41}
\expandafter\def\csname leanline@FamilyClassical11.block_35\endcsname{42}
\expandafter\def\csname leanline@FamilyClassical11.block_36\endcsname{43}
\expandafter\def\csname leanline@FamilyClassical11.block_37\endcsname{44}
\expandafter\def\csname leanline@FamilyClassical11.block_38\endcsname{45}
\expandafter\def\csname leanline@FamilyClassical11.block_39\endcsname{46}
\expandafter\def\csname leanline@FamilyClassical11.block_4\endcsname{11}
\expandafter\def\csname leanline@FamilyClassical11.block_40\endcsname{47}
\expandafter\def\csname leanline@FamilyClassical11.block_41\endcsname{48}
\expandafter\def\csname leanline@FamilyClassical11.block_42\endcsname{49}
\expandafter\def\csname leanline@FamilyClassical11.block_43\endcsname{50}
\expandafter\def\csname leanline@FamilyClassical11.block_44\endcsname{51}
\expandafter\def\csname leanline@FamilyClassical11.block_45\endcsname{52}
\expandafter\def\csname leanline@FamilyClassical11.block_46\endcsname{53}
\expandafter\def\csname leanline@FamilyClassical11.block_47\endcsname{54}
\expandafter\def\csname leanline@FamilyClassical11.block_48\endcsname{55}
\expandafter\def\csname leanline@FamilyClassical11.block_49\endcsname{56}
\expandafter\def\csname leanline@FamilyClassical11.block_5\endcsname{12}
\expandafter\def\csname leanline@FamilyClassical11.block_50\endcsname{57}
\expandafter\def\csname leanline@FamilyClassical11.block_51\endcsname{58}
\expandafter\def\csname leanline@FamilyClassical11.block_52\endcsname{59}
\expandafter\def\csname leanline@FamilyClassical11.block_53\endcsname{60}
\expandafter\def\csname leanline@FamilyClassical11.block_54\endcsname{61}
\expandafter\def\csname leanline@FamilyClassical11.block_55\endcsname{62}
\expandafter\def\csname leanline@FamilyClassical11.block_56\endcsname{63}
\expandafter\def\csname leanline@FamilyClassical11.block_57\endcsname{64}
\expandafter\def\csname leanline@FamilyClassical11.block_58\endcsname{65}
\expandafter\def\csname leanline@FamilyClassical11.block_59\endcsname{66}
\expandafter\def\csname leanline@FamilyClassical11.block_6\endcsname{13}
\expandafter\def\csname leanline@FamilyClassical11.block_60\endcsname{67}
\expandafter\def\csname leanline@FamilyClassical11.block_61\endcsname{68}
\expandafter\def\csname leanline@FamilyClassical11.block_62\endcsname{69}
\expandafter\def\csname leanline@FamilyClassical11.block_63\endcsname{70}
\expandafter\def\csname leanline@FamilyClassical11.block_7\endcsname{14}
\expandafter\def\csname leanline@FamilyClassical11.block_8\endcsname{15}
\expandafter\def\csname leanline@FamilyClassical11.block_9\endcsname{16}
\expandafter\def\csname leanline@FamilyClassical11.encode\endcsname{5}
\expandafter\def\csname leanline@FamilyClassical11.encoded_lower\endcsname{75}
\expandafter\def\csname leanline@FamilyClassical11.goodBlock\endcsname{6}
\expandafter\def\csname leanline@FamilyConnected.Check\endcsname{8}
\expandafter\def\csname leanline@FamilyConnected.E\endcsname{12}
\expandafter\def\csname leanline@FamilyConnected.Reach\endcsname{16}
\expandafter\def\csname leanline@FamilyConnected.Var\endcsname{7}
\expandafter\def\csname leanline@FamilyConnected.Vertex\endcsname{9}
\expandafter\def\csname leanline@FamilyConnected.column_zero\endcsname{43}
\expandafter\def\csname leanline@FamilyConnected.connected\endcsname{48}
\expandafter\def\csname leanline@FamilyConnected.e\endcsname{11}
\expandafter\def\csname leanline@FamilyConnected.nat_shift\endcsname{32}
\expandafter\def\csname leanline@FamilyConnected.same_fiber\endcsname{23}
\expandafter\def\csname leanline@FamilyConnected.shift\endcsname{27}
\expandafter\def\csname leanline@FamilyConnected.vars_connected\endcsname{37}
\expandafter\def\csname leanline@FamilyConnected.via\endcsname{18}
\expandafter\def\csname leanline@FamilyCoverAlternative.classical_lower_twelve\endcsname{10}
\expandafter\def\csname leanline@FamilyGirth.C\endcsname{6}
\expandafter\def\csname leanline@FamilyGirth.V\endcsname{5}
\expandafter\def\csname leanline@FamilyGirth.checks\endcsname{8}
\expandafter\def\csname leanline@FamilyGirth.eight_distinct\endcsname{11}
\expandafter\def\csname leanline@FamilyGirth.eight_edges\endcsname{17}
\expandafter\def\csname leanline@FamilyGirth.exact_eight_certificate\endcsname{29}
\expandafter\def\csname leanline@FamilyGirth.variableVertices\endcsname{9}
\expandafter\def\csname leanline@FamilyNoShort.ColumnPairing\endcsname{18}
\expandafter\def\csname leanline@FamilyNoShort.RowPairing\endcsname{16}
\expandafter\def\csname leanline@FamilyNoShort.base_no_column_pairing\endcsname{23}
\expandafter\def\csname leanline@FamilyNoShort.base_no_four\endcsname{9}
\expandafter\def\csname leanline@FamilyNoShort.base_no_row_pairing\endcsname{22}
\expandafter\def\csname leanline@FamilyNoShort.base_no_six\endcsname{12}
\expandafter\def\csname leanline@FamilyNoShort.e\endcsname{7}
\expandafter\def\csname leanline@FamilyNoShort.no_column_pairing\endcsname{46}
\expandafter\def\csname leanline@FamilyNoShort.no_four\endcsname{26}
\expandafter\def\csname leanline@FamilyNoShort.no_row_pairing\endcsname{38}
\expandafter\def\csname leanline@FamilyNoShort.no_six\endcsname{32}
\expandafter\def\csname leanline@FamilyRing.base\endcsname{10}
\expandafter\def\csname leanline@FamilyRing.conjugateTranspose\endcsname{14}
\expandafter\def\csname leanline@FamilyRing.leftWord\endcsname{19}
\expandafter\def\csname leanline@FamilyRing.mixed_cycle\endcsname{29}
\expandafter\def\csname leanline@FamilyRing.rightWord\endcsname{24}
\expandafter\def\csname leanline@FamilyStopping11.E\endcsname{9}
\expandafter\def\csname leanline@FamilyStopping11.col\endcsname{7}
\expandafter\def\csname leanline@FamilyStopping11.pick\endcsname{14}
\expandafter\def\csname leanline@FamilyStopping11.pick_singleton\endcsname{15}
\expandafter\def\csname leanline@FamilyStopping11.root_0\endcsname{19}
\expandafter\def\csname leanline@FamilyStopping11.root_1\endcsname{20}
\expandafter\def\csname leanline@FamilyStopping11.root_10\endcsname{29}
\expandafter\def\csname leanline@FamilyStopping11.root_11\endcsname{30}
\expandafter\def\csname leanline@FamilyStopping11.root_12\endcsname{31}
\expandafter\def\csname leanline@FamilyStopping11.root_13\endcsname{32}
\expandafter\def\csname leanline@FamilyStopping11.root_14\endcsname{33}
\expandafter\def\csname leanline@FamilyStopping11.root_15\endcsname{34}
\expandafter\def\csname leanline@FamilyStopping11.root_16\endcsname{35}
\expandafter\def\csname leanline@FamilyStopping11.root_17\endcsname{36}
\expandafter\def\csname leanline@FamilyStopping11.root_18\endcsname{37}
\expandafter\def\csname leanline@FamilyStopping11.root_19\endcsname{38}
\expandafter\def\csname leanline@FamilyStopping11.root_2\endcsname{21}
\expandafter\def\csname leanline@FamilyStopping11.root_20\endcsname{39}
\expandafter\def\csname leanline@FamilyStopping11.root_21\endcsname{40}
\expandafter\def\csname leanline@FamilyStopping11.root_22\endcsname{41}
\expandafter\def\csname leanline@FamilyStopping11.root_23\endcsname{42}
\expandafter\def\csname leanline@FamilyStopping11.root_24\endcsname{43}
\expandafter\def\csname leanline@FamilyStopping11.root_25\endcsname{44}
\expandafter\def\csname leanline@FamilyStopping11.root_26\endcsname{45}
\expandafter\def\csname leanline@FamilyStopping11.root_27\endcsname{46}
\expandafter\def\csname leanline@FamilyStopping11.root_28\endcsname{47}
\expandafter\def\csname leanline@FamilyStopping11.root_29\endcsname{48}
\expandafter\def\csname leanline@FamilyStopping11.root_3\endcsname{22}
\expandafter\def\csname leanline@FamilyStopping11.root_30\endcsname{49}
\expandafter\def\csname leanline@FamilyStopping11.root_31\endcsname{50}
\expandafter\def\csname leanline@FamilyStopping11.root_32\endcsname{51}
\expandafter\def\csname leanline@FamilyStopping11.root_33\endcsname{52}
\expandafter\def\csname leanline@FamilyStopping11.root_34\endcsname{53}
\expandafter\def\csname leanline@FamilyStopping11.root_35\endcsname{54}
\expandafter\def\csname leanline@FamilyStopping11.root_36\endcsname{55}
\expandafter\def\csname leanline@FamilyStopping11.root_37\endcsname{56}
\expandafter\def\csname leanline@FamilyStopping11.root_38\endcsname{57}
\expandafter\def\csname leanline@FamilyStopping11.root_39\endcsname{58}
\expandafter\def\csname leanline@FamilyStopping11.root_4\endcsname{23}
\expandafter\def\csname leanline@FamilyStopping11.root_40\endcsname{59}
\expandafter\def\csname leanline@FamilyStopping11.root_41\endcsname{60}
\expandafter\def\csname leanline@FamilyStopping11.root_42\endcsname{61}
\expandafter\def\csname leanline@FamilyStopping11.root_43\endcsname{62}
\expandafter\def\csname leanline@FamilyStopping11.root_5\endcsname{24}
\expandafter\def\csname leanline@FamilyStopping11.root_6\endcsname{25}
\expandafter\def\csname leanline@FamilyStopping11.root_7\endcsname{26}
\expandafter\def\csname leanline@FamilyStopping11.root_8\endcsname{27}
\expandafter\def\csname leanline@FamilyStopping11.root_9\endcsname{28}
\expandafter\def\csname leanline@FamilyStopping11.roots_reject\endcsname{63}
\expandafter\def\csname leanline@FamilyStopping11.row\endcsname{8}
\expandafter\def\csname leanline@FamilyStopping11.row_complete\endcsname{11}
\expandafter\def\csname leanline@FamilyStopping11.stopping_lower_six\endcsname{109}
\expandafter\def\csname leanline@FamilyTranspose.Kernel\endcsname{8}
\expandafter\def\csname leanline@FamilyTranspose.all_constant\endcsname{52}
\expandafter\def\csname leanline@FamilyTranspose.boundary\endcsname{6}
\expandafter\def\csname leanline@FamilyTranspose.constant_of_period\endcsname{7}
\expandafter\def\csname leanline@FamilyTranspose.distance_characterization\endcsname{49}
\expandafter\def\csname leanline@FamilyTranspose.first_constant\endcsname{44}
\expandafter\def\csname leanline@FamilyTranspose.kernel_iff\endcsname{59}
\expandafter\def\csname leanline@FamilyTranspose.nonzero_weight\endcsname{13}
\expandafter\def\csname leanline@FamilyTranspose.pair\endcsname{36}
\expandafter\def\csname leanline@FamilyTranspose.pair_kernel\endcsname{38}
\expandafter\def\csname leanline@FamilyTranspose.pair_nonzero\endcsname{43}
\expandafter\def\csname leanline@FamilyTranspose.second_constant\endcsname{41}
\expandafter\def\csname leanline@FamilyTranspose.second_period\endcsname{25}
\expandafter\def\csname leanline@FamilyTranspose.third_constant\endcsname{16}
\expandafter\def\csname leanline@FamilyTranspose.third_period\endcsname{10}
\expandafter\def\csname leanline@FamilyTranspose.three_even_count\endcsname{6}
\expandafter\def\csname leanline@FamilyTranspose.word\endcsname{11}
\expandafter\def\csname leanline@FamilyTreeStopping.base_stopping_six\endcsname{26}
\expandafter\def\csname leanline@FamilyTreeStopping.no_four\endcsname{8}
\expandafter\def\csname leanline@FamilyTreeStopping.no_six\endcsname{13}
\expandafter\def\csname leanline@FamilyTreeStopping.three_checks\endcsname{18}
\expandafter\def\csname leanline@FamilyWords.Qubit\endcsname{86}
\expandafter\def\csname leanline@FamilyWords.Syndrome\endcsname{28}
\expandafter\def\csname leanline@FamilyWords.U\endcsname{19}
\expandafter\def\csname leanline@FamilyWords.V\endcsname{24}
\expandafter\def\csname leanline@FamilyWords.augmentation\endcsname{41}
\expandafter\def\csname leanline@FamilyWords.augmentation_stabilizer\endcsname{66}
\expandafter\def\csname leanline@FamilyWords.augmentation_witness\endcsname{76}
\expandafter\def\csname leanline@FamilyWords.css_commutation\endcsname{54}
\expandafter\def\csname leanline@FamilyWords.delta\endcsname{12}
\expandafter\def\csname leanline@FamilyWords.delta_add\endcsname{16}
\expandafter\def\csname leanline@FamilyWords.delta_sub\endcsname{14}
\expandafter\def\csname leanline@FamilyWords.exponent\endcsname{11}
\expandafter\def\csname leanline@FamilyWords.stabilizerLeft\endcsname{45}
\expandafter\def\csname leanline@FamilyWords.stabilizerRight\endcsname{49}
\expandafter\def\csname leanline@FamilyWords.sum_shift\endcsname{62}
\expandafter\def\csname leanline@FamilyWords.support\endcsname{95}
\expandafter\def\csname leanline@FamilyWords.supportList\endcsname{90}
\expandafter\def\csname leanline@FamilyWords.witness_kernel\endcsname{34}
\expandafter\def\csname leanline@FamilyWords.witness_nontrivial\endcsname{79}
\expandafter\def\csname leanline@FamilyWords.witness_weight_le\endcsname{106}
\expandafter\def\csname leanline@FamilyWords.word\endcsname{87}
\expandafter\def\csname leanline@FamilyWords.word_zero_outside\endcsname{97}
\expandafter\def\csname leanline@GeneralLP.C\endcsname{8}
\expandafter\def\csname leanline@GeneralLP.HX\endcsname{11}
\expandafter\def\csname leanline@GeneralLP.HZ\endcsname{15}
\expandafter\def\csname leanline@GeneralLP.InSpanX\endcsname{65}
\expandafter\def\csname leanline@GeneralLP.LogicalX\endcsname{69}
\expandafter\def\csname leanline@GeneralLP.P\endcsname{8}
\expandafter\def\csname leanline@GeneralLP.Syndrome\endcsname{19}
\expandafter\def\csname leanline@GeneralLP.XCheck\endcsname{6}
\expandafter\def\csname leanline@GeneralLP.ZCheck\endcsname{7}
\expandafter\def\csname leanline@GeneralLP.augmentation_stabilizer\endcsname{73}
\expandafter\def\csname leanline@GeneralLP.classical_action\endcsname{11}
\expandafter\def\csname leanline@GeneralLP.coefficient_commutation\endcsname{57}
\expandafter\def\csname leanline@GeneralLP.css_commutation\endcsname{37}
\expandafter\def\csname leanline@GeneralLP.full_span_left\endcsname{84}
\expandafter\def\csname leanline@GeneralLP.join\endcsname{9}
\expandafter\def\csname leanline@GeneralLP.nonzero_classical_exists\endcsname{30}
\expandafter\def\csname leanline@GeneralLP.stabilizerLeft\endcsname{24}
\expandafter\def\csname leanline@GeneralLP.stabilizerRight\endcsname{28}
\expandafter\def\csname leanline@GeneralLP.twoFibers\endcsname{17}
\expandafter\def\csname leanline@GeneralLP.twoFibers_kernel\endcsname{19}
\expandafter\def\csname leanline@GeneralLP.twoFibers_nonzero\endcsname{25}
\expandafter\def\csname leanline@GeneralLP.x_span_left\endcsname{43}
\expandafter\def\csname leanline@GeneralLP.x_span_right\endcsname{50}
\expandafter\def\csname leanline@GeneralLP.z_action\endcsname{32}
\expandafter\def\csname leanline@ParametricProjection.ColumnPairing\endcsname{12}
\expandafter\def\csname leanline@ParametricProjection.RowPairing\endcsname{9}
\expandafter\def\csname leanline@ParametricProjection.actual_classical_lower\endcsname{79}
\expandafter\def\csname leanline@ParametricProjection.classical_lower\endcsname{67}
\expandafter\def\csname leanline@ParametricProjection.cyclic_congruence_counterfamily\endcsname{102}
\expandafter\def\csname leanline@ParametricProjection.no_column_pairing\endcsname{56}
\expandafter\def\csname leanline@ParametricProjection.no_four\endcsname{27}
\expandafter\def\csname leanline@ParametricProjection.no_row_pairing\endcsname{45}
\expandafter\def\csname leanline@ParametricProjection.no_six\endcsname{36}
\expandafter\def\csname leanline@ParametricProjection.projected_counterfamily\endcsname{88}
\expandafter\def\csname leanline@ParametricProjection.projected_exponents\endcsname{15}
\expandafter\def\csname leanline@ParametricRing.U\endcsname{18}
\expandafter\def\csname leanline@ParametricRing.V\endcsname{24}
\expandafter\def\csname leanline@ParametricRing.base\endcsname{7}
\expandafter\def\csname leanline@ParametricRing.mixed_cycle\endcsname{28}
\expandafter\def\csname leanline@ParametricRing.star\endcsname{12}
\expandafter\def\csname leanline@ParametricSpecialization.HX_matrix\endcsname{12}
\expandafter\def\csname leanline@ParametricSpecialization.HZ_matrix\endcsname{21}
\expandafter\def\csname leanline@ParametricSpecialization.classical_matrix\endcsname{30}
\expandafter\def\csname leanline@ParametricSpecialization.exponent_matrix\endcsname{7}
\expandafter\def\csname leanline@ParametricSpecialization.left_word\endcsname{34}
\expandafter\def\csname leanline@ParametricSpecialization.logicalX\endcsname{38}
\expandafter\def\csname leanline@ParametricWords.U\endcsname{12}
\expandafter\def\csname leanline@ParametricWords.V\endcsname{18}
\expandafter\def\csname leanline@ParametricWords.actual_logical\endcsname{33}
\expandafter\def\csname leanline@ParametricWords.augmentation_witness\endcsname{29}
\expandafter\def\csname leanline@ParametricWords.e\endcsname{9}
\expandafter\def\csname leanline@ParametricWords.parametric_obstruction\endcsname{76}
\expandafter\def\csname leanline@ParametricWords.source_parametric_obstruction\endcsname{9}
\expandafter\def\csname leanline@ParametricWords.supportList\endcsname{43}
\expandafter\def\csname leanline@ParametricWords.weight_le\endcsname{61}
\expandafter\def\csname leanline@ParametricWords.witness_kernel\endcsname{20}
\expandafter\def\csname leanline@ParametricWords.zero_outside\endcsname{49}
\expandafter\def\csname leanline@SmallStopping.Stopping\endcsname{9}
\expandafter\def\csname leanline@SmallStopping.stopping_extension\endcsname{12}
\expandafter\def\csname leanline@SmallStopping.survivesList\endcsname{32}
\expandafter\def\csname leanline@SmallStopping.survivesList_of_stopping\endcsname{38}
\expandafter\def\csname leanline@Source11.ClassicalHasWeight\endcsname{9}
\expandafter\def\csname leanline@Source11.ClassicalKernel\endcsname{12}
\expandafter\def\csname leanline@Source11.ColumnPairing\endcsname{20}
\expandafter\def\csname leanline@Source11.QuantumHasWeight\endcsname{11}
\expandafter\def\csname leanline@Source11.RowPairing\endcsname{17}
\expandafter\def\csname leanline@Source11.bits\endcsname{10}
\expandafter\def\csname leanline@Source11.classicalDistance\endcsname{28}
\expandafter\def\csname leanline@Source11.classicalRow\endcsname{11}
\expandafter\def\csname leanline@Source11.classicalWitness\endcsname{21}
\expandafter\def\csname leanline@Source11.classical_distance_twelve\endcsname{31}
\expandafter\def\csname leanline@Source11.classical_has_twelve\endcsname{19}
\expandafter\def\csname leanline@Source11.classical_witness_kernel\endcsname{22}
\expandafter\def\csname leanline@Source11.classical_witness_weight\endcsname{26}
\expandafter\def\csname leanline@Source11.css_0\endcsname{6}
\expandafter\def\csname leanline@Source11.css_1\endcsname{9}
\expandafter\def\csname leanline@Source11.css_10\endcsname{36}
\expandafter\def\csname leanline@Source11.css_100\endcsname{306}
\expandafter\def\csname leanline@Source11.css_101\endcsname{309}
\expandafter\def\csname leanline@Source11.css_102\endcsname{312}
\expandafter\def\csname leanline@Source11.css_103\endcsname{315}
\expandafter\def\csname leanline@Source11.css_104\endcsname{318}
\expandafter\def\csname leanline@Source11.css_105\endcsname{321}
\expandafter\def\csname leanline@Source11.css_106\endcsname{324}
\expandafter\def\csname leanline@Source11.css_107\endcsname{327}
\expandafter\def\csname leanline@Source11.css_108\endcsname{330}
\expandafter\def\csname leanline@Source11.css_109\endcsname{333}
\expandafter\def\csname leanline@Source11.css_11\endcsname{39}
\expandafter\def\csname leanline@Source11.css_110\endcsname{336}
\expandafter\def\csname leanline@Source11.css_111\endcsname{339}
\expandafter\def\csname leanline@Source11.css_112\endcsname{342}
\expandafter\def\csname leanline@Source11.css_113\endcsname{345}
\expandafter\def\csname leanline@Source11.css_114\endcsname{348}
\expandafter\def\csname leanline@Source11.css_115\endcsname{351}
\expandafter\def\csname leanline@Source11.css_116\endcsname{354}
\expandafter\def\csname leanline@Source11.css_117\endcsname{357}
\expandafter\def\csname leanline@Source11.css_118\endcsname{360}
\expandafter\def\csname leanline@Source11.css_119\endcsname{363}
\expandafter\def\csname leanline@Source11.css_12\endcsname{42}
\expandafter\def\csname leanline@Source11.css_120\endcsname{366}
\expandafter\def\csname leanline@Source11.css_121\endcsname{369}
\expandafter\def\csname leanline@Source11.css_122\endcsname{372}
\expandafter\def\csname leanline@Source11.css_123\endcsname{375}
\expandafter\def\csname leanline@Source11.css_124\endcsname{378}
\expandafter\def\csname leanline@Source11.css_125\endcsname{381}
\expandafter\def\csname leanline@Source11.css_126\endcsname{384}
\expandafter\def\csname leanline@Source11.css_127\endcsname{387}
\expandafter\def\csname leanline@Source11.css_128\endcsname{390}
\expandafter\def\csname leanline@Source11.css_129\endcsname{393}
\expandafter\def\csname leanline@Source11.css_13\endcsname{45}
\expandafter\def\csname leanline@Source11.css_130\endcsname{396}
\expandafter\def\csname leanline@Source11.css_131\endcsname{399}
\expandafter\def\csname leanline@Source11.css_14\endcsname{48}
\expandafter\def\csname leanline@Source11.css_15\endcsname{51}
\expandafter\def\csname leanline@Source11.css_16\endcsname{54}
\expandafter\def\csname leanline@Source11.css_17\endcsname{57}
\expandafter\def\csname leanline@Source11.css_18\endcsname{60}
\expandafter\def\csname leanline@Source11.css_19\endcsname{63}
\expandafter\def\csname leanline@Source11.css_2\endcsname{12}
\expandafter\def\csname leanline@Source11.css_20\endcsname{66}
\expandafter\def\csname leanline@Source11.css_21\endcsname{69}
\expandafter\def\csname leanline@Source11.css_22\endcsname{72}
\expandafter\def\csname leanline@Source11.css_23\endcsname{75}
\expandafter\def\csname leanline@Source11.css_24\endcsname{78}
\expandafter\def\csname leanline@Source11.css_25\endcsname{81}
\expandafter\def\csname leanline@Source11.css_26\endcsname{84}
\expandafter\def\csname leanline@Source11.css_27\endcsname{87}
\expandafter\def\csname leanline@Source11.css_28\endcsname{90}
\expandafter\def\csname leanline@Source11.css_29\endcsname{93}
\expandafter\def\csname leanline@Source11.css_3\endcsname{15}
\expandafter\def\csname leanline@Source11.css_30\endcsname{96}
\expandafter\def\csname leanline@Source11.css_31\endcsname{99}
\expandafter\def\csname leanline@Source11.css_32\endcsname{102}
\expandafter\def\csname leanline@Source11.css_33\endcsname{105}
\expandafter\def\csname leanline@Source11.css_34\endcsname{108}
\expandafter\def\csname leanline@Source11.css_35\endcsname{111}
\expandafter\def\csname leanline@Source11.css_36\endcsname{114}
\expandafter\def\csname leanline@Source11.css_37\endcsname{117}
\expandafter\def\csname leanline@Source11.css_38\endcsname{120}
\expandafter\def\csname leanline@Source11.css_39\endcsname{123}
\expandafter\def\csname leanline@Source11.css_4\endcsname{18}
\expandafter\def\csname leanline@Source11.css_40\endcsname{126}
\expandafter\def\csname leanline@Source11.css_41\endcsname{129}
\expandafter\def\csname leanline@Source11.css_42\endcsname{132}
\expandafter\def\csname leanline@Source11.css_43\endcsname{135}
\expandafter\def\csname leanline@Source11.css_44\endcsname{138}
\expandafter\def\csname leanline@Source11.css_45\endcsname{141}
\expandafter\def\csname leanline@Source11.css_46\endcsname{144}
\expandafter\def\csname leanline@Source11.css_47\endcsname{147}
\expandafter\def\csname leanline@Source11.css_48\endcsname{150}
\expandafter\def\csname leanline@Source11.css_49\endcsname{153}
\expandafter\def\csname leanline@Source11.css_5\endcsname{21}
\expandafter\def\csname leanline@Source11.css_50\endcsname{156}
\expandafter\def\csname leanline@Source11.css_51\endcsname{159}
\expandafter\def\csname leanline@Source11.css_52\endcsname{162}
\expandafter\def\csname leanline@Source11.css_53\endcsname{165}
\expandafter\def\csname leanline@Source11.css_54\endcsname{168}
\expandafter\def\csname leanline@Source11.css_55\endcsname{171}
\expandafter\def\csname leanline@Source11.css_56\endcsname{174}
\expandafter\def\csname leanline@Source11.css_57\endcsname{177}
\expandafter\def\csname leanline@Source11.css_58\endcsname{180}
\expandafter\def\csname leanline@Source11.css_59\endcsname{183}
\expandafter\def\csname leanline@Source11.css_6\endcsname{24}
\expandafter\def\csname leanline@Source11.css_60\endcsname{186}
\expandafter\def\csname leanline@Source11.css_61\endcsname{189}
\expandafter\def\csname leanline@Source11.css_62\endcsname{192}
\expandafter\def\csname leanline@Source11.css_63\endcsname{195}
\expandafter\def\csname leanline@Source11.css_64\endcsname{198}
\expandafter\def\csname leanline@Source11.css_65\endcsname{201}
\expandafter\def\csname leanline@Source11.css_66\endcsname{204}
\expandafter\def\csname leanline@Source11.css_67\endcsname{207}
\expandafter\def\csname leanline@Source11.css_68\endcsname{210}
\expandafter\def\csname leanline@Source11.css_69\endcsname{213}
\expandafter\def\csname leanline@Source11.css_7\endcsname{27}
\expandafter\def\csname leanline@Source11.css_70\endcsname{216}
\expandafter\def\csname leanline@Source11.css_71\endcsname{219}
\expandafter\def\csname leanline@Source11.css_72\endcsname{222}
\expandafter\def\csname leanline@Source11.css_73\endcsname{225}
\expandafter\def\csname leanline@Source11.css_74\endcsname{228}
\expandafter\def\csname leanline@Source11.css_75\endcsname{231}
\expandafter\def\csname leanline@Source11.css_76\endcsname{234}
\expandafter\def\csname leanline@Source11.css_77\endcsname{237}
\expandafter\def\csname leanline@Source11.css_78\endcsname{240}
\expandafter\def\csname leanline@Source11.css_79\endcsname{243}
\expandafter\def\csname leanline@Source11.css_8\endcsname{30}
\expandafter\def\csname leanline@Source11.css_80\endcsname{246}
\expandafter\def\csname leanline@Source11.css_81\endcsname{249}
\expandafter\def\csname leanline@Source11.css_82\endcsname{252}
\expandafter\def\csname leanline@Source11.css_83\endcsname{255}
\expandafter\def\csname leanline@Source11.css_84\endcsname{258}
\expandafter\def\csname leanline@Source11.css_85\endcsname{261}
\expandafter\def\csname leanline@Source11.css_86\endcsname{264}
\expandafter\def\csname leanline@Source11.css_87\endcsname{267}
\expandafter\def\csname leanline@Source11.css_88\endcsname{270}
\expandafter\def\csname leanline@Source11.css_89\endcsname{273}
\expandafter\def\csname leanline@Source11.css_9\endcsname{33}
\expandafter\def\csname leanline@Source11.css_90\endcsname{276}
\expandafter\def\csname leanline@Source11.css_91\endcsname{279}
\expandafter\def\csname leanline@Source11.css_92\endcsname{282}
\expandafter\def\csname leanline@Source11.css_93\endcsname{285}
\expandafter\def\csname leanline@Source11.css_94\endcsname{288}
\expandafter\def\csname leanline@Source11.css_95\endcsname{291}
\expandafter\def\csname leanline@Source11.css_96\endcsname{294}
\expandafter\def\csname leanline@Source11.css_97\endcsname{297}
\expandafter\def\csname leanline@Source11.css_98\endcsname{300}
\expandafter\def\csname leanline@Source11.css_99\endcsname{303}
\expandafter\def\csname leanline@Source11.css_commutation\endcsname{402}
\expandafter\def\csname leanline@Source11.diagnostic_refuted\endcsname{50}
\expandafter\def\csname leanline@Source11.distance_drop\endcsname{44}
\expandafter\def\csname leanline@Source11.encode_bit_0\endcsname{238}
\expandafter\def\csname leanline@Source11.encode_bit_1\endcsname{242}
\expandafter\def\csname leanline@Source11.encode_bit_10\endcsname{278}
\expandafter\def\csname leanline@Source11.encode_bit_11\endcsname{282}
\expandafter\def\csname leanline@Source11.encode_bit_12\endcsname{286}
\expandafter\def\csname leanline@Source11.encode_bit_13\endcsname{290}
\expandafter\def\csname leanline@Source11.encode_bit_14\endcsname{294}
\expandafter\def\csname leanline@Source11.encode_bit_15\endcsname{298}
\expandafter\def\csname leanline@Source11.encode_bit_16\endcsname{302}
\expandafter\def\csname leanline@Source11.encode_bit_17\endcsname{306}
\expandafter\def\csname leanline@Source11.encode_bit_18\endcsname{310}
\expandafter\def\csname leanline@Source11.encode_bit_19\endcsname{314}
\expandafter\def\csname leanline@Source11.encode_bit_2\endcsname{246}
\expandafter\def\csname leanline@Source11.encode_bit_20\endcsname{318}
\expandafter\def\csname leanline@Source11.encode_bit_21\endcsname{322}
\expandafter\def\csname leanline@Source11.encode_bit_22\endcsname{326}
\expandafter\def\csname leanline@Source11.encode_bit_23\endcsname{330}
\expandafter\def\csname leanline@Source11.encode_bit_24\endcsname{334}
\expandafter\def\csname leanline@Source11.encode_bit_25\endcsname{338}
\expandafter\def\csname leanline@Source11.encode_bit_26\endcsname{342}
\expandafter\def\csname leanline@Source11.encode_bit_27\endcsname{346}
\expandafter\def\csname leanline@Source11.encode_bit_28\endcsname{350}
\expandafter\def\csname leanline@Source11.encode_bit_29\endcsname{354}
\expandafter\def\csname leanline@Source11.encode_bit_3\endcsname{250}
\expandafter\def\csname leanline@Source11.encode_bit_30\endcsname{358}
\expandafter\def\csname leanline@Source11.encode_bit_31\endcsname{362}
\expandafter\def\csname leanline@Source11.encode_bit_32\endcsname{366}
\expandafter\def\csname leanline@Source11.encode_bit_33\endcsname{370}
\expandafter\def\csname leanline@Source11.encode_bit_34\endcsname{374}
\expandafter\def\csname leanline@Source11.encode_bit_35\endcsname{378}
\expandafter\def\csname leanline@Source11.encode_bit_36\endcsname{382}
\expandafter\def\csname leanline@Source11.encode_bit_37\endcsname{386}
\expandafter\def\csname leanline@Source11.encode_bit_38\endcsname{390}
\expandafter\def\csname leanline@Source11.encode_bit_39\endcsname{394}
\expandafter\def\csname leanline@Source11.encode_bit_4\endcsname{254}
\expandafter\def\csname leanline@Source11.encode_bit_40\endcsname{398}
\expandafter\def\csname leanline@Source11.encode_bit_41\endcsname{402}
\expandafter\def\csname leanline@Source11.encode_bit_42\endcsname{406}
\expandafter\def\csname leanline@Source11.encode_bit_43\endcsname{410}
\expandafter\def\csname leanline@Source11.encode_bit_5\endcsname{258}
\expandafter\def\csname leanline@Source11.encode_bit_6\endcsname{262}
\expandafter\def\csname leanline@Source11.encode_bit_7\endcsname{266}
\expandafter\def\csname leanline@Source11.encode_bit_8\endcsname{270}
\expandafter\def\csname leanline@Source11.encode_bit_9\endcsname{274}
\expandafter\def\csname leanline@Source11.exponent\endcsname{9}
\expandafter\def\csname leanline@Source11.exponents\endcsname{8}
\expandafter\def\csname leanline@Source11.freeInput\endcsname{414}
\expandafter\def\csname leanline@Source11.kernel_reconstruct\endcsname{415}
\expandafter\def\csname leanline@Source11.mask_1\endcsname{13}
\expandafter\def\csname leanline@Source11.mask_1024\endcsname{43}
\expandafter\def\csname leanline@Source11.mask_1048576\endcsname{73}
\expandafter\def\csname leanline@Source11.mask_10685880729600\endcsname{190}
\expandafter\def\csname leanline@Source11.mask_10715943403520\endcsname{193}
\expandafter\def\csname leanline@Source11.mask_1073741824\endcsname{103}
\expandafter\def\csname leanline@Source11.mask_10870562226176\endcsname{196}
\expandafter\def\csname leanline@Source11.mask_1099511627776\endcsname{139}
\expandafter\def\csname leanline@Source11.mask_11214161707008\endcsname{199}
\expandafter\def\csname leanline@Source11.mask_11278586216448\endcsname{202}
\expandafter\def\csname leanline@Source11.mask_128\endcsname{34}
\expandafter\def\csname leanline@Source11.mask_131072\endcsname{64}
\expandafter\def\csname leanline@Source11.mask_134217728\endcsname{94}
\expandafter\def\csname leanline@Source11.mask_137438953472\endcsname{124}
\expandafter\def\csname leanline@Source11.mask_14547056328704\endcsname{205}
\expandafter\def\csname leanline@Source11.mask_14559941230592\endcsname{208}
\expandafter\def\csname leanline@Source11.mask_14821932138496\endcsname{211}
\expandafter\def\csname leanline@Source11.mask_15414637625344\endcsname{214}
\expandafter\def\csname leanline@Source11.mask_15719582400512\endcsname{217}
\expandafter\def\csname leanline@Source11.mask_1584842932224\endcsname{142}
\expandafter\def\csname leanline@Source11.mask_15899971026944\endcsname{220}
\expandafter\def\csname leanline@Source11.mask_16\endcsname{25}
\expandafter\def\csname leanline@Source11.mask_16071769718784\endcsname{223}
\expandafter\def\csname leanline@Source11.mask_16204911607808\endcsname{226}
\expandafter\def\csname leanline@Source11.mask_16213503639552\endcsname{229}
\expandafter\def\csname leanline@Source11.mask_16384\endcsname{55}
\expandafter\def\csname leanline@Source11.mask_16777216\endcsname{85}
\expandafter\def\csname leanline@Source11.mask_17029547425792\endcsname{232}
\expandafter\def\csname leanline@Source11.mask_17158394347520\endcsname{235}
\expandafter\def\csname leanline@Source11.mask_17179869184\endcsname{115}
\expandafter\def\csname leanline@Source11.mask_1803888361472\endcsname{145}
\expandafter\def\csname leanline@Source11.mask_1992866922496\endcsname{148}
\expandafter\def\csname leanline@Source11.mask_2\endcsname{16}
\expandafter\def\csname leanline@Source11.mask_2014341758976\endcsname{151}
\expandafter\def\csname leanline@Source11.mask_2048\endcsname{46}
\expandafter\def\csname leanline@Source11.mask_2097152\endcsname{76}
\expandafter\def\csname leanline@Source11.mask_2147483648\endcsname{106}
\expandafter\def\csname leanline@Source11.mask_2199023255552\endcsname{154}
\expandafter\def\csname leanline@Source11.mask_2529735737344\endcsname{157}
\expandafter\def\csname leanline@Source11.mask_256\endcsname{37}
\expandafter\def\csname leanline@Source11.mask_262144\endcsname{67}
\expandafter\def\csname leanline@Source11.mask_2624225017856\endcsname{160}
\expandafter\def\csname leanline@Source11.mask_268435456\endcsname{97}
\expandafter\def\csname leanline@Source11.mask_274877906944\endcsname{127}
\expandafter\def\csname leanline@Source11.mask_304944775168\endcsname{130}
\expandafter\def\csname leanline@Source11.mask_32\endcsname{28}
\expandafter\def\csname leanline@Source11.mask_32768\endcsname{58}
\expandafter\def\csname leanline@Source11.mask_33554432\endcsname{88}
\expandafter\def\csname leanline@Source11.mask_34359738368\endcsname{118}
\expandafter\def\csname leanline@Source11.mask_4\endcsname{19}
\expandafter\def\csname leanline@Source11.mask_4096\endcsname{49}
\expandafter\def\csname leanline@Source11.mask_4114580766720\endcsname{163}
\expandafter\def\csname leanline@Source11.mask_4194304\endcsname{79}
\expandafter\def\csname leanline@Source11.mask_4294967296\endcsname{109}
\expandafter\def\csname leanline@Source11.mask_4316444229632\endcsname{166}
\expandafter\def\csname leanline@Source11.mask_4398046511104\endcsname{169}
\expandafter\def\csname leanline@Source11.mask_4548372463616\endcsname{172}
\expandafter\def\csname leanline@Source11.mask_512\endcsname{40}
\expandafter\def\csname leanline@Source11.mask_524288\endcsname{70}
\expandafter\def\csname leanline@Source11.mask_536870912\endcsname{100}
\expandafter\def\csname leanline@Source11.mask_549755813888\endcsname{133}
\expandafter\def\csname leanline@Source11.mask_64\endcsname{31}
\expandafter\def\csname leanline@Source11.mask_65536\endcsname{61}
\expandafter\def\csname leanline@Source11.mask_67108864\endcsname{91}
\expandafter\def\csname leanline@Source11.mask_68719476736\endcsname{121}
\expandafter\def\csname leanline@Source11.mask_7971461398528\endcsname{175}
\expandafter\def\csname leanline@Source11.mask_8\endcsname{22}
\expandafter\def\csname leanline@Source11.mask_8192\endcsname{52}
\expandafter\def\csname leanline@Source11.mask_8302171783168\endcsname{178}
\expandafter\def\csname leanline@Source11.mask_8370893357056\endcsname{181}
\expandafter\def\csname leanline@Source11.mask_8388608\endcsname{82}
\expandafter\def\csname leanline@Source11.mask_8589934592\endcsname{112}
\expandafter\def\csname leanline@Source11.mask_863290523648\endcsname{136}
\expandafter\def\csname leanline@Source11.mask_8796093022208\endcsname{184}
\expandafter\def\csname leanline@Source11.mask_8800390086656\endcsname{187}
\expandafter\def\csname leanline@Source11.no_column_pairing\endcsname{27}
\expandafter\def\csname leanline@Source11.no_four_cycle\endcsname{31}
\expandafter\def\csname leanline@Source11.no_four_voltage\endcsname{9}
\expandafter\def\csname leanline@Source11.no_row_pairing\endcsname{24}
\expandafter\def\csname leanline@Source11.no_six_cycle\endcsname{37}
\expandafter\def\csname leanline@Source11.no_six_voltage\endcsname{12}
\expandafter\def\csname leanline@Source11.quantumBits\endcsname{7}
\expandafter\def\csname leanline@Source11.quantumDistance\endcsname{29}
\expandafter\def\csname leanline@Source11.quantumWitness\endcsname{6}
\expandafter\def\csname leanline@Source11.quantum_distance_at_most_eleven\endcsname{40}
\expandafter\def\csname leanline@Source11.quantum_has_eleven\endcsname{25}
\expandafter\def\csname leanline@Source11.separates\endcsname{18}
\expandafter\def\csname leanline@Source11.separator\endcsname{7}
\expandafter\def\csname leanline@Source11.separator_kernel\endcsname{13}
\expandafter\def\csname leanline@Source11.voltage\endcsname{7}
\expandafter\def\csname leanline@Source11.witness_kernel\endcsname{8}
\expandafter\def\csname leanline@Source11.witness_nontrivial\endcsname{19}
\expandafter\def\csname leanline@Source11.witness_weight\endcsname{20}
\expandafter\def\csname leanline@Source11.xData\endcsname{14}
\expandafter\def\csname leanline@Source11.xRow\endcsname{8}
\expandafter\def\csname leanline@Source11.x_data_bridge\endcsname{148}
\expandafter\def\csname leanline@Source11.x_row_0\endcsname{16}
\expandafter\def\csname leanline@Source11.x_row_1\endcsname{17}
\expandafter\def\csname leanline@Source11.x_row_10\endcsname{26}
\expandafter\def\csname leanline@Source11.x_row_100\endcsname{116}
\expandafter\def\csname leanline@Source11.x_row_101\endcsname{117}
\expandafter\def\csname leanline@Source11.x_row_102\endcsname{118}
\expandafter\def\csname leanline@Source11.x_row_103\endcsname{119}
\expandafter\def\csname leanline@Source11.x_row_104\endcsname{120}
\expandafter\def\csname leanline@Source11.x_row_105\endcsname{121}
\expandafter\def\csname leanline@Source11.x_row_106\endcsname{122}
\expandafter\def\csname leanline@Source11.x_row_107\endcsname{123}
\expandafter\def\csname leanline@Source11.x_row_108\endcsname{124}
\expandafter\def\csname leanline@Source11.x_row_109\endcsname{125}
\expandafter\def\csname leanline@Source11.x_row_11\endcsname{27}
\expandafter\def\csname leanline@Source11.x_row_110\endcsname{126}
\expandafter\def\csname leanline@Source11.x_row_111\endcsname{127}
\expandafter\def\csname leanline@Source11.x_row_112\endcsname{128}
\expandafter\def\csname leanline@Source11.x_row_113\endcsname{129}
\expandafter\def\csname leanline@Source11.x_row_114\endcsname{130}
\expandafter\def\csname leanline@Source11.x_row_115\endcsname{131}
\expandafter\def\csname leanline@Source11.x_row_116\endcsname{132}
\expandafter\def\csname leanline@Source11.x_row_117\endcsname{133}
\expandafter\def\csname leanline@Source11.x_row_118\endcsname{134}
\expandafter\def\csname leanline@Source11.x_row_119\endcsname{135}
\expandafter\def\csname leanline@Source11.x_row_12\endcsname{28}
\expandafter\def\csname leanline@Source11.x_row_120\endcsname{136}
\expandafter\def\csname leanline@Source11.x_row_121\endcsname{137}
\expandafter\def\csname leanline@Source11.x_row_122\endcsname{138}
\expandafter\def\csname leanline@Source11.x_row_123\endcsname{139}
\expandafter\def\csname leanline@Source11.x_row_124\endcsname{140}
\expandafter\def\csname leanline@Source11.x_row_125\endcsname{141}
\expandafter\def\csname leanline@Source11.x_row_126\endcsname{142}
\expandafter\def\csname leanline@Source11.x_row_127\endcsname{143}
\expandafter\def\csname leanline@Source11.x_row_128\endcsname{144}
\expandafter\def\csname leanline@Source11.x_row_129\endcsname{145}
\expandafter\def\csname leanline@Source11.x_row_13\endcsname{29}
\expandafter\def\csname leanline@Source11.x_row_130\endcsname{146}
\expandafter\def\csname leanline@Source11.x_row_131\endcsname{147}
\expandafter\def\csname leanline@Source11.x_row_14\endcsname{30}
\expandafter\def\csname leanline@Source11.x_row_15\endcsname{31}
\expandafter\def\csname leanline@Source11.x_row_16\endcsname{32}
\expandafter\def\csname leanline@Source11.x_row_17\endcsname{33}
\expandafter\def\csname leanline@Source11.x_row_18\endcsname{34}
\expandafter\def\csname leanline@Source11.x_row_19\endcsname{35}
\expandafter\def\csname leanline@Source11.x_row_2\endcsname{18}
\expandafter\def\csname leanline@Source11.x_row_20\endcsname{36}
\expandafter\def\csname leanline@Source11.x_row_21\endcsname{37}
\expandafter\def\csname leanline@Source11.x_row_22\endcsname{38}
\expandafter\def\csname leanline@Source11.x_row_23\endcsname{39}
\expandafter\def\csname leanline@Source11.x_row_24\endcsname{40}
\expandafter\def\csname leanline@Source11.x_row_25\endcsname{41}
\expandafter\def\csname leanline@Source11.x_row_26\endcsname{42}
\expandafter\def\csname leanline@Source11.x_row_27\endcsname{43}
\expandafter\def\csname leanline@Source11.x_row_28\endcsname{44}
\expandafter\def\csname leanline@Source11.x_row_29\endcsname{45}
\expandafter\def\csname leanline@Source11.x_row_3\endcsname{19}
\expandafter\def\csname leanline@Source11.x_row_30\endcsname{46}
\expandafter\def\csname leanline@Source11.x_row_31\endcsname{47}
\expandafter\def\csname leanline@Source11.x_row_32\endcsname{48}
\expandafter\def\csname leanline@Source11.x_row_33\endcsname{49}
\expandafter\def\csname leanline@Source11.x_row_34\endcsname{50}
\expandafter\def\csname leanline@Source11.x_row_35\endcsname{51}
\expandafter\def\csname leanline@Source11.x_row_36\endcsname{52}
\expandafter\def\csname leanline@Source11.x_row_37\endcsname{53}
\expandafter\def\csname leanline@Source11.x_row_38\endcsname{54}
\expandafter\def\csname leanline@Source11.x_row_39\endcsname{55}
\expandafter\def\csname leanline@Source11.x_row_4\endcsname{20}
\expandafter\def\csname leanline@Source11.x_row_40\endcsname{56}
\expandafter\def\csname leanline@Source11.x_row_41\endcsname{57}
\expandafter\def\csname leanline@Source11.x_row_42\endcsname{58}
\expandafter\def\csname leanline@Source11.x_row_43\endcsname{59}
\expandafter\def\csname leanline@Source11.x_row_44\endcsname{60}
\expandafter\def\csname leanline@Source11.x_row_45\endcsname{61}
\expandafter\def\csname leanline@Source11.x_row_46\endcsname{62}
\expandafter\def\csname leanline@Source11.x_row_47\endcsname{63}
\expandafter\def\csname leanline@Source11.x_row_48\endcsname{64}
\expandafter\def\csname leanline@Source11.x_row_49\endcsname{65}
\expandafter\def\csname leanline@Source11.x_row_5\endcsname{21}
\expandafter\def\csname leanline@Source11.x_row_50\endcsname{66}
\expandafter\def\csname leanline@Source11.x_row_51\endcsname{67}
\expandafter\def\csname leanline@Source11.x_row_52\endcsname{68}
\expandafter\def\csname leanline@Source11.x_row_53\endcsname{69}
\expandafter\def\csname leanline@Source11.x_row_54\endcsname{70}
\expandafter\def\csname leanline@Source11.x_row_55\endcsname{71}
\expandafter\def\csname leanline@Source11.x_row_56\endcsname{72}
\expandafter\def\csname leanline@Source11.x_row_57\endcsname{73}
\expandafter\def\csname leanline@Source11.x_row_58\endcsname{74}
\expandafter\def\csname leanline@Source11.x_row_59\endcsname{75}
\expandafter\def\csname leanline@Source11.x_row_6\endcsname{22}
\expandafter\def\csname leanline@Source11.x_row_60\endcsname{76}
\expandafter\def\csname leanline@Source11.x_row_61\endcsname{77}
\expandafter\def\csname leanline@Source11.x_row_62\endcsname{78}
\expandafter\def\csname leanline@Source11.x_row_63\endcsname{79}
\expandafter\def\csname leanline@Source11.x_row_64\endcsname{80}
\expandafter\def\csname leanline@Source11.x_row_65\endcsname{81}
\expandafter\def\csname leanline@Source11.x_row_66\endcsname{82}
\expandafter\def\csname leanline@Source11.x_row_67\endcsname{83}
\expandafter\def\csname leanline@Source11.x_row_68\endcsname{84}
\expandafter\def\csname leanline@Source11.x_row_69\endcsname{85}
\expandafter\def\csname leanline@Source11.x_row_7\endcsname{23}
\expandafter\def\csname leanline@Source11.x_row_70\endcsname{86}
\expandafter\def\csname leanline@Source11.x_row_71\endcsname{87}
\expandafter\def\csname leanline@Source11.x_row_72\endcsname{88}
\expandafter\def\csname leanline@Source11.x_row_73\endcsname{89}
\expandafter\def\csname leanline@Source11.x_row_74\endcsname{90}
\expandafter\def\csname leanline@Source11.x_row_75\endcsname{91}
\expandafter\def\csname leanline@Source11.x_row_76\endcsname{92}
\expandafter\def\csname leanline@Source11.x_row_77\endcsname{93}
\expandafter\def\csname leanline@Source11.x_row_78\endcsname{94}
\expandafter\def\csname leanline@Source11.x_row_79\endcsname{95}
\expandafter\def\csname leanline@Source11.x_row_8\endcsname{24}
\expandafter\def\csname leanline@Source11.x_row_80\endcsname{96}
\expandafter\def\csname leanline@Source11.x_row_81\endcsname{97}
\expandafter\def\csname leanline@Source11.x_row_82\endcsname{98}
\expandafter\def\csname leanline@Source11.x_row_83\endcsname{99}
\expandafter\def\csname leanline@Source11.x_row_84\endcsname{100}
\expandafter\def\csname leanline@Source11.x_row_85\endcsname{101}
\expandafter\def\csname leanline@Source11.x_row_86\endcsname{102}
\expandafter\def\csname leanline@Source11.x_row_87\endcsname{103}
\expandafter\def\csname leanline@Source11.x_row_88\endcsname{104}
\expandafter\def\csname leanline@Source11.x_row_89\endcsname{105}
\expandafter\def\csname leanline@Source11.x_row_9\endcsname{25}
\expandafter\def\csname leanline@Source11.x_row_90\endcsname{106}
\expandafter\def\csname leanline@Source11.x_row_91\endcsname{107}
\expandafter\def\csname leanline@Source11.x_row_92\endcsname{108}
\expandafter\def\csname leanline@Source11.x_row_93\endcsname{109}
\expandafter\def\csname leanline@Source11.x_row_94\endcsname{110}
\expandafter\def\csname leanline@Source11.x_row_95\endcsname{111}
\expandafter\def\csname leanline@Source11.x_row_96\endcsname{112}
\expandafter\def\csname leanline@Source11.x_row_97\endcsname{113}
\expandafter\def\csname leanline@Source11.x_row_98\endcsname{114}
\expandafter\def\csname leanline@Source11.x_row_99\endcsname{115}
\expandafter\def\csname leanline@Source11.zData\endcsname{15}
\expandafter\def\csname leanline@Source11.zRow\endcsname{11}
\expandafter\def\csname leanline@Source11.z_data_bridge\endcsname{416}
\expandafter\def\csname leanline@Source11.z_row_0\endcsname{284}
\expandafter\def\csname leanline@Source11.z_row_1\endcsname{285}
\expandafter\def\csname leanline@Source11.z_row_10\endcsname{294}
\expandafter\def\csname leanline@Source11.z_row_100\endcsname{384}
\expandafter\def\csname leanline@Source11.z_row_101\endcsname{385}
\expandafter\def\csname leanline@Source11.z_row_102\endcsname{386}
\expandafter\def\csname leanline@Source11.z_row_103\endcsname{387}
\expandafter\def\csname leanline@Source11.z_row_104\endcsname{388}
\expandafter\def\csname leanline@Source11.z_row_105\endcsname{389}
\expandafter\def\csname leanline@Source11.z_row_106\endcsname{390}
\expandafter\def\csname leanline@Source11.z_row_107\endcsname{391}
\expandafter\def\csname leanline@Source11.z_row_108\endcsname{392}
\expandafter\def\csname leanline@Source11.z_row_109\endcsname{393}
\expandafter\def\csname leanline@Source11.z_row_11\endcsname{295}
\expandafter\def\csname leanline@Source11.z_row_110\endcsname{394}
\expandafter\def\csname leanline@Source11.z_row_111\endcsname{395}
\expandafter\def\csname leanline@Source11.z_row_112\endcsname{396}
\expandafter\def\csname leanline@Source11.z_row_113\endcsname{397}
\expandafter\def\csname leanline@Source11.z_row_114\endcsname{398}
\expandafter\def\csname leanline@Source11.z_row_115\endcsname{399}
\expandafter\def\csname leanline@Source11.z_row_116\endcsname{400}
\expandafter\def\csname leanline@Source11.z_row_117\endcsname{401}
\expandafter\def\csname leanline@Source11.z_row_118\endcsname{402}
\expandafter\def\csname leanline@Source11.z_row_119\endcsname{403}
\expandafter\def\csname leanline@Source11.z_row_12\endcsname{296}
\expandafter\def\csname leanline@Source11.z_row_120\endcsname{404}
\expandafter\def\csname leanline@Source11.z_row_121\endcsname{405}
\expandafter\def\csname leanline@Source11.z_row_122\endcsname{406}
\expandafter\def\csname leanline@Source11.z_row_123\endcsname{407}
\expandafter\def\csname leanline@Source11.z_row_124\endcsname{408}
\expandafter\def\csname leanline@Source11.z_row_125\endcsname{409}
\expandafter\def\csname leanline@Source11.z_row_126\endcsname{410}
\expandafter\def\csname leanline@Source11.z_row_127\endcsname{411}
\expandafter\def\csname leanline@Source11.z_row_128\endcsname{412}
\expandafter\def\csname leanline@Source11.z_row_129\endcsname{413}
\expandafter\def\csname leanline@Source11.z_row_13\endcsname{297}
\expandafter\def\csname leanline@Source11.z_row_130\endcsname{414}
\expandafter\def\csname leanline@Source11.z_row_131\endcsname{415}
\expandafter\def\csname leanline@Source11.z_row_14\endcsname{298}
\expandafter\def\csname leanline@Source11.z_row_15\endcsname{299}
\expandafter\def\csname leanline@Source11.z_row_16\endcsname{300}
\expandafter\def\csname leanline@Source11.z_row_17\endcsname{301}
\expandafter\def\csname leanline@Source11.z_row_18\endcsname{302}
\expandafter\def\csname leanline@Source11.z_row_19\endcsname{303}
\expandafter\def\csname leanline@Source11.z_row_2\endcsname{286}
\expandafter\def\csname leanline@Source11.z_row_20\endcsname{304}
\expandafter\def\csname leanline@Source11.z_row_21\endcsname{305}
\expandafter\def\csname leanline@Source11.z_row_22\endcsname{306}
\expandafter\def\csname leanline@Source11.z_row_23\endcsname{307}
\expandafter\def\csname leanline@Source11.z_row_24\endcsname{308}
\expandafter\def\csname leanline@Source11.z_row_25\endcsname{309}
\expandafter\def\csname leanline@Source11.z_row_26\endcsname{310}
\expandafter\def\csname leanline@Source11.z_row_27\endcsname{311}
\expandafter\def\csname leanline@Source11.z_row_28\endcsname{312}
\expandafter\def\csname leanline@Source11.z_row_29\endcsname{313}
\expandafter\def\csname leanline@Source11.z_row_3\endcsname{287}
\expandafter\def\csname leanline@Source11.z_row_30\endcsname{314}
\expandafter\def\csname leanline@Source11.z_row_31\endcsname{315}
\expandafter\def\csname leanline@Source11.z_row_32\endcsname{316}
\expandafter\def\csname leanline@Source11.z_row_33\endcsname{317}
\expandafter\def\csname leanline@Source11.z_row_34\endcsname{318}
\expandafter\def\csname leanline@Source11.z_row_35\endcsname{319}
\expandafter\def\csname leanline@Source11.z_row_36\endcsname{320}
\expandafter\def\csname leanline@Source11.z_row_37\endcsname{321}
\expandafter\def\csname leanline@Source11.z_row_38\endcsname{322}
\expandafter\def\csname leanline@Source11.z_row_39\endcsname{323}
\expandafter\def\csname leanline@Source11.z_row_4\endcsname{288}
\expandafter\def\csname leanline@Source11.z_row_40\endcsname{324}
\expandafter\def\csname leanline@Source11.z_row_41\endcsname{325}
\expandafter\def\csname leanline@Source11.z_row_42\endcsname{326}
\expandafter\def\csname leanline@Source11.z_row_43\endcsname{327}
\expandafter\def\csname leanline@Source11.z_row_44\endcsname{328}
\expandafter\def\csname leanline@Source11.z_row_45\endcsname{329}
\expandafter\def\csname leanline@Source11.z_row_46\endcsname{330}
\expandafter\def\csname leanline@Source11.z_row_47\endcsname{331}
\expandafter\def\csname leanline@Source11.z_row_48\endcsname{332}
\expandafter\def\csname leanline@Source11.z_row_49\endcsname{333}
\expandafter\def\csname leanline@Source11.z_row_5\endcsname{289}
\expandafter\def\csname leanline@Source11.z_row_50\endcsname{334}
\expandafter\def\csname leanline@Source11.z_row_51\endcsname{335}
\expandafter\def\csname leanline@Source11.z_row_52\endcsname{336}
\expandafter\def\csname leanline@Source11.z_row_53\endcsname{337}
\expandafter\def\csname leanline@Source11.z_row_54\endcsname{338}
\expandafter\def\csname leanline@Source11.z_row_55\endcsname{339}
\expandafter\def\csname leanline@Source11.z_row_56\endcsname{340}
\expandafter\def\csname leanline@Source11.z_row_57\endcsname{341}
\expandafter\def\csname leanline@Source11.z_row_58\endcsname{342}
\expandafter\def\csname leanline@Source11.z_row_59\endcsname{343}
\expandafter\def\csname leanline@Source11.z_row_6\endcsname{290}
\expandafter\def\csname leanline@Source11.z_row_60\endcsname{344}
\expandafter\def\csname leanline@Source11.z_row_61\endcsname{345}
\expandafter\def\csname leanline@Source11.z_row_62\endcsname{346}
\expandafter\def\csname leanline@Source11.z_row_63\endcsname{347}
\expandafter\def\csname leanline@Source11.z_row_64\endcsname{348}
\expandafter\def\csname leanline@Source11.z_row_65\endcsname{349}
\expandafter\def\csname leanline@Source11.z_row_66\endcsname{350}
\expandafter\def\csname leanline@Source11.z_row_67\endcsname{351}
\expandafter\def\csname leanline@Source11.z_row_68\endcsname{352}
\expandafter\def\csname leanline@Source11.z_row_69\endcsname{353}
\expandafter\def\csname leanline@Source11.z_row_7\endcsname{291}
\expandafter\def\csname leanline@Source11.z_row_70\endcsname{354}
\expandafter\def\csname leanline@Source11.z_row_71\endcsname{355}
\expandafter\def\csname leanline@Source11.z_row_72\endcsname{356}
\expandafter\def\csname leanline@Source11.z_row_73\endcsname{357}
\expandafter\def\csname leanline@Source11.z_row_74\endcsname{358}
\expandafter\def\csname leanline@Source11.z_row_75\endcsname{359}
\expandafter\def\csname leanline@Source11.z_row_76\endcsname{360}
\expandafter\def\csname leanline@Source11.z_row_77\endcsname{361}
\expandafter\def\csname leanline@Source11.z_row_78\endcsname{362}
\expandafter\def\csname leanline@Source11.z_row_79\endcsname{363}
\expandafter\def\csname leanline@Source11.z_row_8\endcsname{292}
\expandafter\def\csname leanline@Source11.z_row_80\endcsname{364}
\expandafter\def\csname leanline@Source11.z_row_81\endcsname{365}
\expandafter\def\csname leanline@Source11.z_row_82\endcsname{366}
\expandafter\def\csname leanline@Source11.z_row_83\endcsname{367}
\expandafter\def\csname leanline@Source11.z_row_84\endcsname{368}
\expandafter\def\csname leanline@Source11.z_row_85\endcsname{369}
\expandafter\def\csname leanline@Source11.z_row_86\endcsname{370}
\expandafter\def\csname leanline@Source11.z_row_87\endcsname{371}
\expandafter\def\csname leanline@Source11.z_row_88\endcsname{372}
\expandafter\def\csname leanline@Source11.z_row_89\endcsname{373}
\expandafter\def\csname leanline@Source11.z_row_9\endcsname{293}
\expandafter\def\csname leanline@Source11.z_row_90\endcsname{374}
\expandafter\def\csname leanline@Source11.z_row_91\endcsname{375}
\expandafter\def\csname leanline@Source11.z_row_92\endcsname{376}
\expandafter\def\csname leanline@Source11.z_row_93\endcsname{377}
\expandafter\def\csname leanline@Source11.z_row_94\endcsname{378}
\expandafter\def\csname leanline@Source11.z_row_95\endcsname{379}
\expandafter\def\csname leanline@Source11.z_row_96\endcsname{380}
\expandafter\def\csname leanline@Source11.z_row_97\endcsname{381}
\expandafter\def\csname leanline@Source11.z_row_98\endcsname{382}
\expandafter\def\csname leanline@Source11.z_row_99\endcsname{383}
\expandafter\def\csname leanline@SourceDistance.C\endcsname{8}
\expandafter\def\csname leanline@SourceDistance.ClassicalHasWeight\endcsname{46}
\expandafter\def\csname leanline@SourceDistance.LogicalX\endcsname{17}
\expandafter\def\csname leanline@SourceDistance.LogicalZ\endcsname{21}
\expandafter\def\csname leanline@SourceDistance.QuantumHasWeight\endcsname{50}
\expandafter\def\csname leanline@SourceDistance.c_action\endcsname{11}
\expandafter\def\csname leanline@SourceDistance.classical_weight_iff\endcsname{56}
\expandafter\def\csname leanline@SourceDistance.logicalX_iff\endcsname{25}
\expandafter\def\csname leanline@SourceDistance.logicalZ_iff\endcsname{35}
\expandafter\def\csname leanline@SourceDistance.quantum_weight_iff\endcsname{60}
\expandafter\def\csname leanline@SourceDistance.source_counterfamily\endcsname{67}
\expandafter\def\csname leanline@SourceModel.HX\endcsname{13}
\expandafter\def\csname leanline@SourceModel.HZ\endcsname{17}
\expandafter\def\csname leanline@SourceModel.InSpanX\endcsname{35}
\expandafter\def\csname leanline@SourceModel.InSpanZ\endcsname{39}
\expandafter\def\csname leanline@SourceModel.P\endcsname{11}
\expandafter\def\csname leanline@SourceModel.P_action\endcsname{25}
\expandafter\def\csname leanline@SourceModel.XCheck\endcsname{8}
\expandafter\def\csname leanline@SourceModel.ZCheck\endcsname{9}
\expandafter\def\csname leanline@SourceModel.inSpanX_iff\endcsname{43}
\expandafter\def\csname leanline@SourceModel.inSpanZ_iff\endcsname{60}
\expandafter\def\csname leanline@SourceModel.join\endcsname{21}
\expandafter\def\csname leanline@SourceModel.join_surjective\endcsname{18}
\expandafter\def\csname leanline@SourceModel.qubit_card\endcsname{7}
\expandafter\def\csname leanline@SourceModel.transpose\endcsname{7}
\expandafter\def\csname leanline@SourceModel.x_action\endcsname{9}
\expandafter\def\csname leanline@SourceModel.x_span_left\endcsname{44}
\expandafter\def\csname leanline@SourceModel.x_span_right\endcsname{51}
\expandafter\def\csname leanline@SourceModel.xcheck_card\endcsname{11}
\expandafter\def\csname leanline@SourceModel.z_action\endcsname{34}
\expandafter\def\csname leanline@SourceModel.z_span_left\endcsname{19}
\expandafter\def\csname leanline@SourceModel.z_span_right\endcsname{27}
\expandafter\def\csname leanline@SourceModel.zcheck_card\endcsname{14}
\expandafter\def\csname leanline@SourceSymmetry.XHasWeight\endcsname{29}
\expandafter\def\csname leanline@SourceSymmetry.ZHasWeight\endcsname{34}
\expandafter\def\csname leanline@SourceSymmetry.join_transpose\endcsname{12}
\expandafter\def\csname leanline@SourceSymmetry.sector_weights\endcsname{40}
\expandafter\def\csname leanline@SourceSymmetry.swapQ\endcsname{8}
\expandafter\def\csname leanline@SourceSymmetry.transpose_twice\endcsname{24}
\expandafter\def\csname leanline@SourceSymmetry.transpose_weight\endcsname{18}
\expandafter\def\csname leanline@SourceWordBridge.InSpan\endcsname{11}
\expandafter\def\csname leanline@SourceWordBridge.Kernel\endcsname{8}
\expandafter\def\csname leanline@SourceWordBridge.QuantumHasWeight\endcsname{14}
\expandafter\def\csname leanline@SourceWordBridge.block_weight_iff\endcsname{19}
\expandafter\def\csname leanline@SourceWordBridge.classicalDistance\endcsname{42}
\expandafter\def\csname leanline@SourceWordBridge.classical_distance_eq\endcsname{49}
\expandafter\def\csname leanline@SourceWordBridge.classical_exists\endcsname{37}
\expandafter\def\csname leanline@SourceWordBridge.css_commutation\endcsname{63}
\expandafter\def\csname leanline@SourceWordBridge.quantumDistance\endcsname{41}
\expandafter\def\csname leanline@SourceWordBridge.quantum_distance_eq\endcsname{44}
\expandafter\def\csname leanline@SourceWordBridge.quantum_exists\endcsname{33}
\expandafter\def\csname leanline@SourceWordBridge.quantum_weight_iff\endcsname{27}
\expandafter\def\csname leanline@SourceWordBridge.source_distance_drop\endcsname{55}
\expandafter\def\csname leanline@TannerSix.NoFour\endcsname{6}
\expandafter\def\csname leanline@TannerSix.NoSix\endcsname{9}
\expandafter\def\csname leanline@TannerSix.ThreeChecks\endcsname{13}
\expandafter\def\csname leanline@TannerSix.partner\endcsname{16}
\expandafter\def\csname leanline@TannerSix.stopping_lower_six\endcsname{47}
\expandafter\def\csname leanline@TannerSix.two_more\endcsname{32}
\expandafter\def\csname leanline@VoltageCycles.Check\endcsname{7}
\expandafter\def\csname leanline@VoltageCycles.Edge\endcsname{11}
\expandafter\def\csname leanline@VoltageCycles.NoFourVoltage\endcsname{32}
\expandafter\def\csname leanline@VoltageCycles.NoSixVoltage\endcsname{35}
\expandafter\def\csname leanline@VoltageCycles.Variable\endcsname{8}
\expandafter\def\csname leanline@VoltageCycles.check_base_ne\endcsname{14}
\expandafter\def\csname leanline@VoltageCycles.no_four\endcsname{40}
\expandafter\def\csname leanline@VoltageCycles.no_six\endcsname{48}
\expandafter\def\csname leanline@VoltageCycles.variable_base_ne\endcsname{22}

\newcommand{\lean}[2]{\href{\proofbase/#1.lean\#L\csname leanline@#2\endcsname}{\nolinkurl{#2}}}
\newcommand{\proofsource}[4]{\href{\proofbase/#1.lean\#L#2-L#3}{\nolinkurl{#4}}}
\newcommand{\leanlisting}[1]{\csname leanlisting@#1\endcsname}
% Generated complete-declaration source listings.
\expandafter\def\csname leanlisting@source-objects\endcsname{\par\smallskip\noindent\textbf{Checked source:} \lean{GeneralLP}{GeneralLP.HX}; \lean{GeneralLP}{GeneralLP.HZ}; \lean{GeneralLP}{GeneralLP.Syndrome}; \lean{GeneralLP}{GeneralLP.stabilizerLeft}; \lean{GeneralLP}{GeneralLP.stabilizerRight}; \lean{GeneralLP}{GeneralLP.InSpanX}; \lean{GeneralLP}{GeneralLP.LogicalX}.\par\smallskip\lstinputlisting{listings/source-objects.tex}}
\expandafter\def\csname leanlisting@source-span\endcsname{\par\smallskip\noindent\textbf{Checked source:} \lean{SourceSectors}{SourceModel.InSpanX}; \lean{SourceSectors}{SourceModel.inSpanX_iff}.\par\smallskip\lstinputlisting{listings/source-span.tex}}
\expandafter\def\csname leanlisting@source-minimum\endcsname{\par\smallskip\noindent\textbf{Checked source:} \lean{SourceWordBridge}{SourceWordBridge.Kernel}; \lean{SourceWordBridge}{SourceWordBridge.InSpan}; \lean{SourceWordBridge}{SourceWordBridge.QuantumHasWeight}.\par\smallskip\lstinputlisting{listings/source-minimum.tex}}
\expandafter\def\csname leanlisting@parametric-word\endcsname{\par\smallskip\noindent\textbf{Checked source:} \lean{ParametricWords}{ParametricWords.e}; \lean{ParametricWords}{ParametricWords.U}; \lean{ParametricWords}{ParametricWords.V}; \lean{ParametricWords}{ParametricWords.witness_kernel}.\par\smallskip\lstinputlisting{listings/parametric-word.tex}}
\expandafter\def\csname leanlisting@separator-span\endcsname{\par\smallskip\noindent\textbf{Checked source:} \lean{GeneralLP}{GeneralLP.augmentation_stabilizer}; \lean{GeneralLP}{GeneralLP.full_span_left}.\par\smallskip\lstinputlisting{listings/separator-span.tex}}
\expandafter\def\csname leanlisting@logical-separation\endcsname{\par\smallskip\noindent\textbf{Checked source:} \lean{ParametricWords}{ParametricWords.augmentation_witness}; \lean{ParametricWords}{ParametricWords.actual_logical}.\par\smallskip\lstinputlisting{listings/logical-separation.tex}}
\expandafter\def\csname leanlisting@physical-weight\endcsname{\par\smallskip\noindent\textbf{Checked source:} \lean{ParametricWords}{ParametricWords.supportList}; \lean{ParametricWords}{ParametricWords.zero_outside}; \lean{ParametricWords}{ParametricWords.weight_le}.\par\smallskip\lstinputlisting{listings/physical-weight.tex}}
\expandafter\def\csname leanlisting@source-obstruction\endcsname{\par\smallskip\noindent\textbf{Checked source:} \lean{ParametricSource}{ParametricWords.source_parametric_obstruction}.\par\smallskip\lstinputlisting{listings/source-obstruction.tex}}
\expandafter\def\csname leanlisting@cancelled-fibers\endcsname{\par\smallskip\noindent\textbf{Checked source:} \lean{CoverWeight}{CoverWeight.push}; \lean{CoverWeight}{CoverWeight.zero_push_fiber_lower}; \lean{CoverWeight}{CoverWeight.zero_push_weight_doubles}.\par\smallskip\lstinputlisting{listings/cancelled-fibers.tex}}
\expandafter\def\csname leanlisting@stopping-projection\endcsname{\par\smallskip\noindent\textbf{Checked source:} \lean{CyclicCover}{CyclicCover.projected_support_stopping}.\par\smallskip\lstinputlisting{listings/stopping-projection.tex}}
\expandafter\def\csname leanlisting@cover-bound\endcsname{\par\smallskip\noindent\textbf{Checked source:} \lean{CoverDistance}{CoverDistance.lower_bound}.\par\smallskip\lstinputlisting{listings/cover-bound.tex}}
\expandafter\def\csname leanlisting@finite-certificate\endcsname{\par\smallskip\noindent\textbf{Checked source:} \lean{FamilyDistance11}{FamilyClassical11.encoded_lower}; \lean{FamilyDistance11}{FamilyClassical11.actual_classical_lower}.\par\smallskip\lstinputlisting{listings/finite-certificate.tex}}
\expandafter\def\csname leanlisting@finite-weight-bridge\endcsname{\par\smallskip\noindent\textbf{Checked source:} \lean{FamilyBaseDistance}{FamilyBaseGraph.packWord_weight}; \lean{FamilyBaseDistance}{FamilyBaseGraph.code_lower_twelve}.\par\smallskip\lstinputlisting{listings/finite-weight-bridge.tex}}
\expandafter\def\csname leanlisting@tree-stopping\endcsname{\par\smallskip\noindent\textbf{Checked source:} \lean{TannerSix}{TannerSix.stopping_lower_six}.\par\smallskip\lstinputlisting{listings/tree-stopping.tex}}
\expandafter\def\csname leanlisting@cover-application\endcsname{\par\smallskip\noindent\textbf{Checked source:} \lean{FamilyCoverAlternative}{FamilyCoverAlternative.classical_lower_twelve}.\par\smallskip\lstinputlisting{listings/cover-application.tex}}
\expandafter\def\csname leanlisting@classical-upper\endcsname{\par\smallskip\noindent\textbf{Checked source:} \lean{ClassicalUpper20}{ClassicalUpper20.word}; \lean{ClassicalUpper20}{ClassicalUpper20.kernel}; \lean{ClassicalUpper20}{ClassicalUpper20.nonzero}; \lean{ClassicalUpper20}{ClassicalUpper20.classical_distance_upper}.\par\smallskip\lstinputlisting{listings/classical-upper.tex}}
\expandafter\def\csname leanlisting@transpose-period\endcsname{\par\smallskip\noindent\textbf{Checked source:} \lean{FamilyTranspose}{FamilyTranspose.boundary}; \lean{FamilyTranspose}{FamilyTranspose.Kernel}; \lean{FamilyTranspose}{FamilyTranspose.third_period}.\par\smallskip\lstinputlisting{listings/transpose-period.tex}}
\expandafter\def\csname leanlisting@transpose-constant\endcsname{\par\smallskip\noindent\textbf{Checked source:} \lean{TransposeConstant}{FamilyTranspose.kernel_iff}.\par\smallskip\lstinputlisting{listings/transpose-constant.tex}}
\expandafter\def\csname leanlisting@transpose-distance\endcsname{\par\smallskip\noindent\textbf{Checked source:} \lean{TransposeDistance}{FamilyTranspose.nonzero_weight}; \lean{TransposeDistance}{FamilyTranspose.distance_characterization}.\par\smallskip\lstinputlisting{listings/transpose-distance.tex}}
\expandafter\def\csname leanlisting@girth-eight\endcsname{\par\smallskip\noindent\textbf{Checked source:} \lean{FamilyGirth}{FamilyGirth.eight_distinct}; \lean{FamilyGirth}{FamilyGirth.eight_edges}; \lean{FamilyGirth}{FamilyGirth.exact_eight_certificate}.\par\smallskip\lstinputlisting{listings/girth-eight.tex}}
\expandafter\def\csname leanlisting@connected\endcsname{\par\smallskip\noindent\textbf{Checked source:} \lean{FamilyConnected}{FamilyConnected.shift}; \lean{FamilyConnected}{FamilyConnected.connected}.\par\smallskip\lstinputlisting{listings/connected.tex}}
\expandafter\def\csname leanlisting@arithmetic-family\endcsname{\par\smallskip\noindent\textbf{Checked source:} \lean{ParametricProjection}{ParametricProjection.projected_counterfamily}.\par\smallskip\lstinputlisting{listings/arithmetic-family.tex}}
\expandafter\def\csname leanlisting@minimum-bridge\endcsname{\par\smallskip\noindent\textbf{Checked source:} \lean{SourceWordBridge}{SourceWordBridge.block_weight_iff}; \lean{SourceWordBridge}{SourceWordBridge.quantum_weight_iff}; \lean{SourceWordBridge}{SourceWordBridge.quantum_distance_eq}; \lean{SourceWordBridge}{SourceWordBridge.classical_distance_eq}.\par\smallskip\lstinputlisting{listings/minimum-bridge.tex}}
\expandafter\def\csname leanlisting@final-distance\endcsname{\par\smallskip\noindent\textbf{Checked source:} \lean{SourceWordBridge}{SourceWordBridge.source_distance_drop}.\par\smallskip\lstinputlisting{listings/final-distance.tex}}

\newcolumntype{P}[1]{>{\raggedright\arraybackslash}p{#1}}
\lstset{basicstyle=\small\ttfamily,columns=fullflexible,keepspaces=true,
 frame=single,framerule=.3pt,breaklines=true,showstringspaces=false,
 escapeinside={(*@}{@*)},
 literate={∀}{{$\forall$}}1 {∃}{{$\exists$}}1 {≠}{{$\ne$}}1
 {∧}{{$\land$}}1 {→}{{$\to$}}1 {≤}{{$\le$}}1 {¬}{{$\neg$}}1
 {∑}{{$\sum$}}1 {↔}{{$\leftrightarrow$}}1 {∣}{{$\mid$}}1
 {⟨}{{$\langle$}}1 {⟩}{{$\rangle$}}1}
\title{Distance loss in girth-eight cyclic lifted-product codes\\
An infinite counterfamily with checked Lean proofs}
\author{}
\date{25 October 2026}
\begin{document}
\maketitle
\begin{abstract}
We construct a connected cyclic lifted-product family whose classical Tanner
graph has girth exactly eight and whose exponent matrix satisfies neither
full row nor full column partition constraint, but whose quantum distance
is strictly smaller than its classical distance. For every positive multiple
$L$ of eleven, the same fully populated $3\times4$ exponent matrix gives
$25L$ physical qubits and
\[
 d_Q\leq 11<12\leq d_C\leq20.
\]
The transpose classical code has distance $2L$, so requiring its distance
to be at least $d_C$ does not remove the obstruction. The logical operator
is pure $X$, with support in both product blocks; a fiber-sum functional
separates it from the entire stabilizer span. A two-parameter template
gives such a logical of weight at most eleven for every positive cyclic
length, including even lengths and parameters with cancelling terms.
Classical lower bounds pass from the eleven-fold base through a parity
projection and a stopping-set bound. The results refute a proposed
girth/partition sufficiency test motivated by Raveendran, Declercq and
Vasi\'c's distance-equality question; their general sufficient-condition
question remains open here. Every proof in the manuscript has a checked
Lean counterpart, including the source-matrix and full-span translations.
\end{abstract}

\section{The result and the source question}\label{sec:answer}

Consider the fixed integer exponent matrix
\begin{equation}\label{eq:E}
 E=\begin{pmatrix}0&0&0&0\\0&1&5&3\\0&2&10&7\end{pmatrix}.
\end{equation}
At cyclic length $L$, each entry is reduced modulo $L$ and specifies a
circulant permutation, so an exponent zero is an identity block. Let $C$
be its $3L\times4L$ binary classical check matrix. We use the symmetric
single-base lifted product of Raveendran, Declercq and Vasi\'c
\cite[Eq.~(9)]{RDV}, with precise conventions in Section~\ref{sec:model}.

\begin{theorem}[Connected counterfamily]\label{thm:main}
For every positive integer $L$ divisible by eleven, the construction from
\eqref{eq:E} has the following properties:
\begin{enumerate}
\item The classical Tanner graph is connected and has girth exactly eight.
\item Neither full row nor full column RCPC of
\cite[Definitions 2--5]{RDV} holds.
\item Its classical and quantum distances satisfy
$d_Q\leq11<12\leq d_C\leq20$.
\item Every nonzero word in $\ker C^{\mathsf T}$ has weight $2L$, and
such a word exists. In particular, the transpose-code distance is at least $d_C$.
\end{enumerate}
\end{theorem}

The proof is assembled in Sections~\ref{sec:logical}--\ref{sec:geometry}.
Its decisive features are an explicit logical in the full stabilizer
quotient and a uniform classical lower bound. No exact value of $d_Q$ for
the whole family is asserted. The bounded distances also mean that this
is a finite-length obstruction family, rather than an asymptotically good
code family.

Raveendran, Declercq and Vasi\'c's preprint, first posted on 10 March 2025,
asks in Section VI for sufficient constraints under which the quantum
distance equals the distance of the underlying classical code. Their
partition constraints identify a low-weight obstruction at the stabilizer
generator scale. The question motivates testing whether classical girth
and exclusion of those full partitions already suffice. Theorem~\ref{thm:main}
answers that test negatively. The test is our proposed criterion; it is
not a conjecture attributed to the source authors.

\begin{longtable}{@{}P{.26\textwidth}P{.29\textwidth}P{.39\textwidth}@{}}
\toprule Question&Source and precise locator&Answer in this paper\\\midrule\endhead
Which constraints guarantee $d_Q=d_C$?&
Raveendran--Declercq--Vasi\'c, arXiv:2503.07567v1,
10 March 2025, Section VI; construction Eq.~(9).&
Still unresolved in general. The counterfamily rules out the particular
tests in the next two rows.\\[4pt]
Do classical girth at least eight and absence of full RCPC suffice?&
Our diagnostic follow-up to that source question; RCPC as in
Definitions 2--5 of the same version.&
No. Theorem~\ref{thm:main} gives an infinite family with
$d_Q\leq11<12\leq d_C$.\\[4pt]
Does additionally excluding a lighter transpose code suffice?&
Our strengthened diagnostic follow-up.&
No. The same family has transpose distance $2L\geq d_C$
(Proposition~\ref{prop:transpose}).\\
\bottomrule
\end{longtable}

The source's Theorem 2 gives an RCPC implication; its Section V-B and
conclusion also describe a recipe excluding stabilizer-weight logicals.
Our weight-eleven obstruction does not refute that scale-seven exclusion
for a $3\times4$ base. Instead, it shows why removing that obstruction
need not reach the larger classical distance.

\subsection*{Relation to earlier graph questions}
Freij-Hollanti, Morris and \v Sapokait\.{e}~\cite{FHM} study the Tanner
graphs of this single-base construction, prove sector graph isomorphism
and connectivity results, and ask in Section 5 for sufficient girth-eight
conditions and their consequences. Earlier project work answers the
girth criterion using rectangle voltages. That local graph information
does not determine the minimum of the stabilizer quotient. Here the
classical graph passes the short-cycle tests and is connected, while a
logical supported in both product blocks still reduces distance.
The connection is therefore mathematical, but the invariant and the
answered diagnostic are different.

The stopping ingredient is classical. Kelley, Sridhara, Xu and
Rosenthal~\cite[Theorem III.1 and the following note]{Kelley} give the
tree bound and its stopping-set interpretation. Their graph-cover
description also explains why the occupied image of a lifted codeword
is a stopping set. We provide checked versions of the degree-three,
girth-eight case and the parity-projection argument needed here; we
do not claim these classical principles as new.

Two more recent distance results address different implications.
Bhardwaj et al.~\cite[Appendix G, Theorem 8]{Bhardwaj} give an
upper bound with a full binary row-rank premise; this does not supply
the lower bound required for $d_Q=d_C$. Hirasaki and
Lee~\cite[Corollary 6]{Hirasaki} compare quantum distances before and
after a compatible group-extension lift. Their odd-index upper bound
and their reverse bound with an additional unchanged-logical-dimension
premise compare quantum codes with quantum codes, rather than the
quantum and classical distances in the present question. Neither
result is used as a premise of our counterfamily proof.

For the central sources, the chronology to the manuscript date is:
\begin{center}\small
\begin{tabular}{@{}P{.25\textwidth}P{.22\textwidth}P{.23\textwidth}P{.22\textwidth}@{}}
\toprule Source&Earliest relevant date&Inspected milestone&Elapsed to 25 Oct. 2026\\\midrule
Raveendran et al.~\cite{RDV}&10 March 2025&arXiv v1, same date&594 days\\
Freij-Hollanti et al.~\cite{FHM}&10 June 2026&arXiv v1, same date&137 days\\
Kelley et al.~\cite{Kelley}&2004 (year precision)&ISIT 2004 proceedings&About 22 years\\
Bhardwaj et al.~\cite{Bhardwaj}&30 July 2026&v1; internal date 24 Aug. 2026&87 days; 62 from internal date\\
Hirasaki--Lee~\cite{Hirasaki}&30 July 2026&v1; internal date 24 Aug. 2026&87 days; 62 from internal date\\
\bottomrule
\end{tabular}
\end{center}
These are source ages, not measurements of how long a problem has been
unsolved. Appendix~\ref{app:references} explains the date convention and
the bounded literature, GitHub and X audit.

\section{The actual source matrices and distances}\label{sec:model}
Write $G_L=\Z/L\Z$ and
$R_L=\F[x,x^{-1}]/(x^L-1)$. For $e\in G_L$ put
\[
 P^e(t,s)=\begin{cases}1&t=s+e,\\0&\text{otherwise},\end{cases}
 \qquad t,s\in G_L.
\]
Binary expansion replaces $x^e$ by $P^e$. With
$B_{ij}=x^{E_{ij}}$, set $B^*=\overline B^{\mathsf T}$, where
$\overline{x^e}=x^{-e}$. The matrices are
\begin{equation}\label{eq:source}
 C=\operatorname{expand}(B),\quad
 \HX=\operatorname{expand}[B\otimes I_4\mid I_3\otimes B^*],\quad
 \HZ=\operatorname{expand}[I_4\otimes B\mid B^*\otimes I_3].
\end{equation}
The Kronecker products are taken in the ring before expansion. Thus both
quantum check matrices have $12L$ rows and $25L$ columns: there is one
cyclic fiber.

Number rows and columns from zero. A quantum word has blocks
$u(j,k,s)$, $j,k<4$, and $v(l,i,s)$, $l,i<3$, represented by matrices
$U\in R_L^{4\times4}$, $V\in R_L^{3\times3}$. All weights count
nonzero binary coefficients across both blocks. The $Z$ check at
$(j,i,t)$ evaluates to
\begin{equation}\label{eq:syndrome}
 \sum_{k<4}u(j,k,t-E_{ik})+
 \sum_{l<3}v(l,i,t+E_{lj}).
\end{equation}
Consequently the pure-$X$ kernel equation is
\begin{equation}\label{eq:ringkernel}
 UB^{\mathsf T}=\overline B^{\mathsf T}V.
\end{equation}
For arbitrary binary check coefficients $w(i,j,s)$, the full $X$ stabilizer
span consists exactly of the pairs
\begin{align}
 u(i,j,s)&=\sum_{k<3}w(k,j,s+E_{ki}),\label{eq:spanleft}\\
 v(i,j,s)&=\sum_{k<4}w(i,k,s-E_{jk}).\label{eq:spanright}
\end{align}
Here the index $i$ in \eqref{eq:spanleft} ranges below four; in
\eqref{eq:spanright} it ranges below three. These formulas quantify
over all lifted-check coefficients.

\begin{proposition}[Matrix and quotient bridges]\label{prop:bridges}
The matrices \eqref{eq:source} satisfy $\HZ\HX^{\mathsf T}=0$.
The $\HZ$ check action and full $\HX$ row span are exactly
\eqref{eq:syndrome}--\eqref{eq:spanright}. For the fixed single-base
family, transposing both square qubit blocks exchanges the two CSS
logical sectors and preserves coefficient weight.
\end{proposition}
\begin{proof}
Expand each monomial using $t=s+e$. Summing first over the single
nonzero circulant entry gives \eqref{eq:syndrome} and the two span
formulas. Substitution of \eqref{eq:spanleft}--\eqref{eq:spanright}
into \eqref{eq:syndrome} cancels the two double sums, proving
commutation for arbitrary check coefficients and hence for every pair
of check rows. Transposing the two square coordinate blocks and the
check type exchanges the $X$ and $Z$ coefficient matrices. This is a
coordinate permutation, so it preserves weight and transports the
entire spans. The cyclic coordinate is unchanged.
The checked matrix identities and equivalences are
\lean{SourceSectors}{SourceModel.inSpanX_iff},
\lean{SourceSectors}{SourceModel.inSpanZ_iff},
\lean{SourceWordBridge}{SourceWordBridge.css_commutation} and
\lean{SourceSymmetry}{SourceSymmetry.sector_weights}.
\end{proof}

Let $d_C$ be the minimum weight of a nonzero word in $\ker C$.
For $d_Q$ use the minimum of the two pure CSS logical spectra,
\[
 (\ker\HZ\setminus\row\HX)\ \cup\
 (\ker\HX\setminus\row\HZ),
\]
in the physical binary coordinates. The formal definitions use minima
over all words, with explicit existence witnesses; the quantum quotient
is nonzero. Our bound also applies directly to usual Pauli distance,
because its exhibited operator is a nontrivial pure-$X$ Pauli operator.
The translation of the actual source minima is checked in
\lean{SourceWordBridge}{SourceWordBridge.quantum_distance_eq} and
\lean{SourceWordBridge}{SourceWordBridge.classical_distance_eq}.

\section{A persistent logical supported in both blocks}\label{sec:logical}

The obstruction is more transparent in two parameters. For $a,b\in G_L$
define
\begin{equation}\label{eq:param}
 E(a,b)=\begin{pmatrix}0&0&0&0\\0&a&b&3a\\0&2a&2b&b+2a\end{pmatrix}.
\end{equation}
The fixed family is $a=1,b=5$. For this template put
\begin{equation}\label{eq:UV}
 U=\begin{pmatrix}
 1&0&0&0\\0&x^{-2a}&0&0\\0&0&x^{-2b}&0\\
 x^{-2a}&x^{-4a}+x^{-2a}&x^{-(b+2a)}&x^{-4a}
 \end{pmatrix},\qquad
 V=\begin{pmatrix}0&0&1\\0&1&0\\1&0&0\end{pmatrix}.
\end{equation}

\begin{theorem}[Universal template obstruction]\label{thm:template}
For every positive $L$ and every $a,b\in G_L$, the pair \eqref{eq:UV}
is in $\ker\HZ\setminus\row\HX$ for the actual matrices of
\eqref{eq:param}, and has weight at most eleven.
\end{theorem}
\begin{proof}
Direct multiplication gives \eqref{eq:ringkernel}. The first three
rows on both sides are, respectively,
\[
 (1,1,1),\quad (x^{-2a},x^{-a},1),\quad(x^{-2b},x^{-b},1).
\]
The last rows agree after cancellation in characteristic two; their
three entries are
\[
 x^{-(b+2a)},\qquad x^{-3a},\qquad 1.
\]
These identities hold with formal monomials and therefore also after
any cyclic quotient. They are checked in
\lean{ParametricRing}{ParametricRing.mixed_cycle} and, directly in
coefficients, \lean{ParametricWords}{ParametricWords.witness_kernel}.

To exclude every stabilizer combination, define
\begin{equation}\label{eq:ell}
 \ell(u)=\sum_{s\in G_L}u(0,0,s)+\sum_{s\in G_L}u(1,0,s)\in\F.
\end{equation}
For a span word \eqref{eq:spanleft}, each of the two sums is
$\sum_{k,s}w(k,0,s)$, since translation permutes $G_L$. They cancel.
For \eqref{eq:UV} the first sum is one and the second zero, so
$\ell(U)=1$. Thus the word lies outside the full $X$ row span.
Finally, $U$ contains eight displayed monomial terms and $V$ three.
Their coefficient support is contained in a list of length eleven.
Coinciding terms may cancel, and cannot increase the weight. The
source-matched conclusion, including matrix commutation, is
\lean{ParametricSource}{ParametricWords.source_parametric_obstruction}.
\end{proof}

The word is pure $X$ throughout. Its use of both product blocks is
distinct from combining $X$ and $Z$ Pauli components. This distinction
matters: a search confined to one product block would miss the obstruction.

For the fixed exponents, the row-pairing equality
$E_{0c}+E_{2c}=2E_{1c}$ holds in the first three columns but fails in the
fourth. The last row of $U$ compensates for that failure. Absence of a
\emph{full} partition therefore does not remove all partial-pairing
logical mechanisms.

\section{The classical lower bound}\label{sec:classical}

For an exponent array $e$ over a finite cyclic group $A$, write a classical
word as $z(j,s)$, with boundary
\[
 (\partial_e z)(i,t)=\sum_{j<4}z(j,t-e_{ij}).
\]
A stopping set means a nonempty variable set for which every check has
either zero or at least two selected neighbors.

\begin{lemma}[Cover projection bound]\label{lem:cover}
Let $p:A\to B$ be an additive map of finite cyclic groups and
$e'=p\circ e$. Suppose every nonzero word of $\ker\partial_{e'}$
has weight at least $d$, and every stopping set of its Tanner incidence
has size at least $s$. Then every nonzero word of $\ker\partial_e$
has weight at least $\min(d,2s)$.
\end{lemma}
\begin{proof}
For $z\ne0$ in the upstairs kernel, define its binary parity projection
and its occupied image by
\[
 y(j,r)=\sum_{t:p(t)=r}z(j,t),\qquad
 T=\{(j,p(t)):z(j,t)\ne0\}.
\]
Projection commutes with the boundary, so $\partial_{e'}y=0$.
Also $\wt(y)\leq\wt(z)$. If $y\ne0$, the code bound gives
$\wt(z)\geq d$.

If $y=0$, each occupied fiber contributes a positive even number of
ones, hence $\wt(z)\geq2|T|$. The set $T$ is a stopping set:
choose any lifted selected neighbor of a touched projected check.
The lifted kernel equation supplies a second selected neighbor at
that check. A monomial row has only one neighbor of each column type,
so the two neighbors have different column types and remain distinct
after projection. Thus $|T|\geq s$, proving the second bound.
This argument is \lean{CoverDistance}{CoverDistance.lower_bound};
\lean{CyclicCover}{CyclicCover.projected_support_stopping} and
\lean{CoverWeight}{CoverWeight.zero_push_weight_doubles} are its two
essential cancellation bridges.
\end{proof}

The use of $T$ as well as $y$ is necessary. A parity projection may
vanish even though the occupied image is nonempty, especially in even
covers. The proof assumes neither odd cover degree nor unchanged rank.

\begin{lemma}[The eleven-fold base]\label{lem:base}
For \eqref{eq:E} at $L=11$, every nonzero classical kernel word has
weight at least twelve, and every stopping set has size at least six.
\end{lemma}
\begin{proof}
For the code bound, the finite certificate reconstructs every
$44$-bit source-kernel word as an encoding of thirteen free bits.
All $2^{13}=8192$ inputs are checked: the encoded word is zero or its
population count is at least twelve. Reconstruction is proved from
the actual source rows; it is not an assumption that a displayed
list spans the kernel. The bit packing, dot-product and population-count
bridges identify this count with coefficient weight.
The checked chain is
\lean{FamilySource11}{FamilyBase11.kernel_reconstruct},
\lean{FamilyDistance11}{FamilyClassical11.encoded_lower},
\lean{FamilyDistance11}{FamilyClassical11.actual_classical_lower}, and
\lean{FamilyBaseDistance}{FamilyBaseGraph.code_lower_twelve}.

For the stopping bound, the base has variable degree three and no
four- or six-cycle, as checked by its finite voltage tests. Choose a
variable $v$ in a stopping set. Its three check neighbors have selected
partners $w_0,w_1,w_2$, distinct by exclusion of four-cycles. At $w_0$,
choose two other checks and their selected partners $x,y$.
Four-cycle exclusion prevents $x=y$ or either partner being $v$;
six-cycle exclusion prevents them being $w_1,w_2$. Together with
$w_0$ being distinct, this yields the six distinct variables
$v,w_0,w_1,w_2,x,y$. This is the degree-three, girth-eight Tanner-tree
bound credited above. Its general checked version is
\lean{TannerSix}{TannerSix.stopping_lower_six}; its actual base
incidence application is
\lean{FamilyTreeStopping}{FamilyTreeStopping.base_stopping_six}.
\end{proof}

\begin{corollary}[Uniform classical lower bound]\label{cor:lower}
If $L>0$ and $11\mid L$, the classical code of \eqref{eq:E} has
$d_C\geq12$.
\end{corollary}
\begin{proof}
Use reduction $G_L\to G_{11}$ in Lemma~\ref{lem:cover}, with
$d=12$ and $s=6$ from Lemma~\ref{lem:base}. The exponent array reduces
to the eleven-fold base and $\min(12,2\cdot6)=12$.
The source-action bridge is
\lean{SourceDistance}{SourceDistance.c_action}; the application is
\lean{FamilyCoverAlternative}{FamilyCoverAlternative.classical_lower_twelve}.
\end{proof}

\begin{corollary}[Arithmetic extensions]\label{cor:param}
Let $L>0$, and suppose an additive map $p:G_L\to G_{11}$ sends
$a$ to $1$ and $b$ to $5$. Then the template \eqref{eq:param} has
no classical four- or six-cycle, has neither full RCPC, and satisfies
$d_Q\leq11<12\leq d_C$. In particular this holds when $11\mid L$
and $a,b$ reduce to $1,5$ modulo eleven.
\end{corollary}
\begin{proof}
The exponent array projects to \eqref{eq:E} modulo eleven. The same
cover bound proves the classical inequality, and
Theorem~\ref{thm:template} supplies the logical. A four-cycle, six-cycle
or full partition upstairs would project to a forbidden voltage or
partition equality at eleven; the base exclusions are specified in
Section~\ref{sec:geometry}. The complete arithmetic implication is
\lean{ParametricProjection}{ParametricProjection.projected_counterfamily},
with the stated congruence specialization in
\lean{ParametricProjection}{ParametricProjection.cyclic_congruence_counterfamily}.
\end{proof}
We assert connectedness, the bound $d_C\leq20$ and transpose distance
$2L$ only for the fixed specialization of Theorem~\ref{thm:main}.

\section{Upper bound, transpose code and graph geometry}\label{sec:geometry}

\begin{proposition}[A classical word of weight at most twenty]\label{prop:upper}
For the fixed family with $L>0$ and $11\mid L$, $d_C\leq20$.
\end{proposition}
\begin{proof}
The following column vector is a polynomial kernel word for $B$:
\begin{equation}\label{eq:twenty}
 \begin{pmatrix}
 x^4+x^5+x^7+x^8+x^9+x^{11}\\
 x^2+x^3+x^6+x^7+x^8+x^{11}\\
 1+x^2+x^3+x^6\\
 1+x^4+x^5+x^9
 \end{pmatrix}.
\end{equation}
Multiplying by each row of $B$ cancels all terms in characteristic two.
There are twenty displayed monomials, so the coefficient weight is at
most twenty after reduction. The third component is nonzero: its
exponents $0,2,3,6$ remain distinct after reduction modulo eleven,
and hence cannot cancel upstairs. The coefficient kernel identity,
nonzero bridge and support bound are proved in
\lean{ClassicalUpper20}{ClassicalUpper20.classical_distance_upper}.
\end{proof}

\begin{proposition}[Exact transpose distance]\label{prop:transpose}
For the fixed matrix \eqref{eq:E} and every positive $L$, the words
of $\ker C^{\mathsf T}$ are exactly the triples of constant fiber
functions $(z_0,z_1,z_2)$ with $z_0+z_1+z_2=0$.
Every nonzero such word has weight $2L$, and a nonzero such word exists.
\end{proposition}
\begin{proof}
The actual transpose equations, with their plus sign, are
\[
 f_j(t):=\sum_{i<3}z_i(t+E_{ij})=0\quad(j<4).
\]
There is a direct characteristic-two identity
\begin{equation}\label{eq:bezout}
 \sum_{j=0}^3\sum_{r\in A_j} f_j(t+r)=z_2(t+1)+z_2(t+2),
\end{equation}
where
\[
\begin{split}
 A_0&=\{1,4,6,7,8,9\},\quad A_1=\{0,2,4,8,9\},\\
 A_2&=\{3,4\},\qquad\quad A_3=\{0,1,2,3,4,6,7\}.
\end{split}
\]
Thus $z_2$ has period one and is constant. The equations $f_0=f_1=0$
then make $z_1$ constant; $f_0=0$ makes $z_0$ constant as well and
gives their sum zero. Conversely such a constant triple satisfies
every equation. A nonzero even binary triple has exactly two ones,
giving weight $2L$; $(1,1,0)$ supplies existence.
Identity~\eqref{eq:bezout} is
\lean{FamilyTranspose}{FamilyTranspose.third_period}; the characterization
and count are
\lean{TransposeConstant}{FamilyTranspose.kernel_iff} and
\lean{TransposeDistance}{FamilyTranspose.distance_characterization}.
The actual transpose action is checked in
\lean{FamilyBounds}{DistanceFamily.transpose_c_action}.
\end{proof}

For the comparison $d_C\leq2L$ at every positive $L$, set two entire
classical variable fibers to one and the other two to zero. Each
classical check then sees two ones. This word has weight $2L$;
\lean{FamilyBounds}{DistanceFamily.classical_distance_upper} checks the
kernel, nonzero and weight claims. Consequently no nonzero transpose
word is lighter than $d_C$, as assembled in
\lean{FamilyBounds}{DistanceFamily.stronger_counterfamily}.

\begin{proposition}[Classical graph and full partitions]\label{prop:graph}
For $L>0$ divisible by eleven, the classical graph of \eqref{eq:E}
is connected, has girth exactly eight and has neither full RCPC.
\end{proposition}
\begin{proof}
A classical incidence joins variable $(j,s)$ to check
$(i,s+E_{ij})$. A simple four-cycle has two distinct row types
and two distinct column types and forces
\[
 E_{ic}-E_{jc}+E_{jd}-E_{id}=0.
\]
A simple six-cycle has three distinct row and column types and forces
\[
 E_{ic}-E_{jc}+E_{jd}-E_{kd}+E_{kf}-E_{if}=0.
\]
Distinctness of the types follows from the unique neighbor of any
specified type at a monomial incidence; the graph/cycle bridge is
\lean{VoltageCycles}{VoltageCycles.no_four} and
\lean{VoltageCycles}{VoltageCycles.no_six}.
The finite tests at eleven show that neither voltage can vanish
with the indicated distinct types. Any equality upstairs reduces
to one at eleven, so both cycle lengths are excluded for every
$11\mid L$.

There is an explicit eight-cycle, written alternately with check
vertices $c$ and variable vertices $v$:
\[
 c_{0,0},\ v_{0,0},\ c_{1,0},\ v_{1,-1},\ c_{2,1},\ v_{0,1},\
 c_{1,1},\ v_{1,0},\ c_{0,0}.
\]
Each displayed edge satisfies the incidence equation. The four
variables and four checks are separately distinct since $L\geq11$.
The graph is simple and bipartite, so its girth is exactly eight.
The voltage exclusions and distinctness are
\lean{FamilyGirth}{FamilyGirth.exact_eight_certificate}; all eight
edges are checked in \lean{FamilyGirth}{FamilyGirth.eight_edges}.

For three rows a full row RCPC means that distinct $i,j,k$ satisfy
$E_{ic}+E_{jc}=2E_{kc}$ in every column, including the odd-size
self-pair convention. For four columns a full column RCPC partitions
the columns into two pairs with equal sums in every row. The finite
tests exclude both at eleven; again projection excludes them
upstairs. These are
\lean{FamilyNoShort}{FamilyNoShort.no_row_pairing} and
\lean{FamilyNoShort}{FamilyNoShort.no_column_pairing}.

Finally the zero row connects all variable types at a common fiber,
and the zero column connects the check types. The entry $E_{11}=1$
gives a walk from variable $(0,s)$ to $(0,s+1)$ through variable
$(1,s)$ and check $(1,s+1)$. Iterating reaches every cyclic fiber,
so all vertices are connected. This argument works for every
positive $L$ and is \lean{FamilyConnected}{FamilyConnected.connected}.
\end{proof}

\begin{proof}[Proof of Theorem~\ref{thm:main}]
Proposition~\ref{prop:graph} proves the graph and RCPC assertions.
Theorem~\ref{thm:template}, specialized to $a=1,b=5$, gives
$d_Q\leq11$. Corollary~\ref{cor:lower} and
Proposition~\ref{prop:upper} give $12\leq d_C\leq20$.
Proposition~\ref{prop:transpose} and the two-fiber classical word
give the transpose comparison. The specialization of all source
matrices and full logical predicates is checked in
\lean{ParametricSpecialization}{ParametricSpecialization.logicalX}.
The source-distance conclusion is
\lean{SourceWordBridge}{SourceWordBridge.source_distance_drop}.
\end{proof}

\section{Formal coverage and reproducibility}\label{sec:formal}
The proofs use Lean 4.36.0-rc1 and mathlib commit
\texttt{5f7b4198d22a9c82e1c698a379334fe128095c82}.
The admitted counterfamily closure contains fifty authored modules.
All were built from their recorded sources and rebuilt in an isolated
directory with original research compiled objects excluded. Pinned
cached mathlib, dependency and standard-library objects were reused;
this was not a fresh mathlib rebuild. The axiom audit admits only
\texttt{propext}, \texttt{Classical.choice} and \texttt{Quot.sound}.
Native-decision prototypes and external search exclusions are outside
this closure.

Table~\ref{tab:coverage} covers each manuscript proof, including the
translations and applications. Appendices~\ref{app:lean-map}--\ref{app:lean-assembly}
explain the formal objects and give complete central Lean proofs, from
the stabilizer separator and cover bound to the source-distance conclusion.
Appendix~\ref{app:lean-links} gives current and immutable URLs for all fifty
admitted modules. Every theorem link in the text targets its declaration
at commit \texttt{98eb91cd9049}, whose source bytes match the checked
builds. The companion \href{PROOF-GUIDE.html}{proof guide} and source
archive provide hashes, logs and editable files in addition to the PDF
exposition. Generated appendix listings are exact source extracts with
typographical substitutions for Unicode symbols, not new proof modules.

\begin{longtable}{@{}P{.27\textwidth}P{.40\textwidth}P{.27\textwidth}@{}}
\caption{Complete proof-to-Lean correspondence. Full names, source hashes,
imports and checked logs are linked in the companion guide.}\label{tab:coverage}\\
\toprule Manuscript proof&Faithful formal result&Objects and scope\\\midrule\endfirsthead
\toprule Manuscript proof&Faithful formal result&Objects and scope\\\midrule\endhead
Proposition~\ref{prop:bridges}; source minima&
\lean{SourceSectors}{SourceModel.inSpanX_iff};
\lean{SourceSectors}{SourceModel.inSpanZ_iff};
\lean{SourceWordBridge}{SourceWordBridge.quantum_weight_iff};
\lean{SourceWordBridge}{SourceWordBridge.css_commutation};
\lean{SourceSymmetry}{SourceSymmetry.sector_weights}&
Actual coefficient matrices, arbitrary binary row coefficients, both
sectors, all physical words and weights.\\[4pt]
Theorem~\ref{thm:template}&
\lean{ParametricRing}{ParametricRing.mixed_cycle};
\lean{ParametricSource}{ParametricWords.source_parametric_obstruction}&
Every positive $L$, arbitrary $a,b\in G_L$; full row span, cancellation
and physical support bound.\\[4pt]
Lemma~\ref{lem:cover}&
\lean{CoverDistance}{CoverDistance.lower_bound};
\lean{CyclicCover}{CyclicCover.kernel_push};
\lean{CyclicCover}{CyclicCover.projected_support_stopping};
\lean{CoverWeight}{CoverWeight.zero_push_weight_doubles}&
Additive projection; nonzero upstairs word; parity and occupied-support
cases treated separately.\\[4pt]
Lemma~\ref{lem:base}, code certificate&
\lean{FamilySource11}{FamilyBase11.kernel_reconstruct};
\lean{FamilyDistance11}{FamilyClassical11.actual_classical_lower};
\lean{FamilyBaseDistance}{FamilyBaseGraph.code_lower_twelve}&
Every source-kernel word; all thirteen-bit inputs; incidence, dot-product
and Hamming-weight bridges.\\[4pt]
Lemma~\ref{lem:base}, stopping bound&
\lean{TannerSix}{TannerSix.stopping_lower_six};
\lean{FamilyTreeStopping}{FamilyTreeStopping.base_stopping_six}&
Nonempty stopping sets; degree at least three and no four/six cycles;
actual base incidence.\\[4pt]
Corollary~\ref{cor:lower}&
\lean{FamilyCoverAlternative}{FamilyCoverAlternative.classical_lower_twelve};
\lean{SourceDistance}{SourceDistance.c_action}&
All $L>0$, $11\mid L$; actual classical source words.\\[4pt]
Corollary~\ref{cor:param}&
\lean{ParametricProjection}{ParametricProjection.projected_counterfamily};
\lean{ParametricProjection}{ParametricProjection.cyclic_congruence_counterfamily}&
Explicit additive map sending $a,b$ to $1,5$; no implicit unit,
connectedness or rank premise.\\[4pt]
Proposition~\ref{prop:upper}&
\lean{ClassicalUpper20}{ClassicalUpper20.kernel};
\lean{ClassicalUpper20}{ClassicalUpper20.nonzero};
\lean{ClassicalUpper20}{ClassicalUpper20.classical_distance_upper}&
Polynomial word, nonzero under reduction modulo eleven, at most twenty
physical coordinates.\\[4pt]
Proposition~\ref{prop:transpose}; two-fiber comparison&
\lean{FamilyTranspose}{FamilyTranspose.third_period};
\lean{TransposeConstant}{FamilyTranspose.kernel_iff};
\lean{TransposeDistance}{FamilyTranspose.distance_characterization};
\lean{FamilyBounds}{DistanceFamily.stronger_counterfamily}&
Actual $C^{\mathsf T}$, every positive $L$, all kernel words and nonempty
minimum; comparison with actual $d_C$.\\[4pt]
Proposition~\ref{prop:graph}&
\lean{VoltageCycles}{VoltageCycles.no_four};
\lean{VoltageCycles}{VoltageCycles.no_six};
\lean{FamilyNoShort}{FamilyNoShort.no_row_pairing};
\lean{FamilyNoShort}{FamilyNoShort.no_column_pairing};
\lean{FamilyGirth}{FamilyGirth.exact_eight_certificate};
\lean{FamilyGirth}{FamilyGirth.eight_edges};
\lean{FamilyConnected}{FamilyConnected.connected}&
Simple-cycle bridges, checked base exclusions, explicit eight-cycle
vertices and edges, actual Tanner reachability.\\[4pt]
Theorem~\ref{thm:main}, assembly&
\lean{FamilyTheorem}{DistanceFamily.infinite_counterfamily};
\lean{SourceWordBridge}{SourceWordBridge.source_distance_drop};
\lean{ParametricSpecialization}{ParametricSpecialization.logicalX};
preceding bound/graph theorems&
Same source construction and minima, $L>0$, $11\mid L$; no exact
quantum distance added.\\
\bottomrule
\end{longtable}

The formal appendix follows the mathematical argument in dependency order:
actual source objects and full spans; the logical obstruction; parity
projection and the base certificate; graph and transpose properties; and
the final all-word distance statement. The dependency diagram in
Appendix~\ref{app:lean-map} separates these branches and their trust
boundaries. External distance searches are not premises of the proofs.

\section{Implications and the next bounded question}\label{sec:limits}
The counterfamily separates local graph and full-partition exclusions
from quantum distance preservation. Connectedness rules out explaining
it as disjoint copies of the eleven-fold example, and the exact
transpose characterization rules out a lighter transpose code as the
cause. What remains is a logical using both product blocks, produced
by a partial pairing that full RCPC tests do not detect.

The source's general equality problem is not answered by this negative
test. A bounded follow-up is to prove or refute
\begin{equation}\label{eq:next}
 d_Q\geq\min(d_C,12)
\end{equation}
for normalized, fully populated $3\times4$ cyclic exponent matrices
whose classical Tanner graph has girth at least eight, which have no
full RCPC, whose transpose distance is at least $d_C$, and which also
obey the following partial-pairing exclusion: for every choice of
distinct rows $i,j,k$, no value among
\[
 E_{ic}+E_{jc}-2E_{kc},\qquad c=0,1,2,3,
\]
occurs three or more times. Normalization means the first row and
first column are zero; the entry tests are stated in this convention.
The fixed counterfamily violates this additional test.

This is a project-derived follow-up, proposed after the counterexample
and the bounded diagnostic recorded in the supporting evidence. It
is not a proved criterion or a new author-posed conjecture. The
remaining obstruction is control of all small logical supports in
the full quotient, especially weights eight through eleven. Two
separately checked generic lemmas provide structural progress, but
neither completes the source-support classification or proves
\eqref{eq:next}. A complete exact-distance certificate for a selected
finite instance is a separate alternative. Verifying the source's
scale-seven recipe is a source-claim audit, rather than a newly open
mathematical problem.

\noindent\textbf{Question answered.} The proposed girth/RCPC criterion,
even with connectedness and the transpose-distance comparison, is
false by Theorem~\ref{thm:main}.\\
\textbf{Question still open.} General sufficient constraints guaranteeing
$d_Q=d_C$ in the source construction.\\
\textbf{Next proposed target (follow-up).} Prove or refute the truncated
bound \eqref{eq:next} under the explicit strengthened hypotheses above.

\appendix
\clearpage
\section{The checked Lean argument: objects and dependency map}\label{app:lean-map}

This appendix presents the formal argument inside the paper. It explains how
the source objects are represented and prints the complete central declaration
bodies used in the distance obstruction, together with the applications needed
to recover the mathematical theorems. Appendix~\ref{app:lean-links} then gives
current and immutable links for every admitted source file. These links
supplement the displayed proofs; they are also the route to routine coordinate
lemmas and the large finite row certificates not printed in full.

All listings below are generated from the checked sources at commit
\texttt{98eb91cd9049}. They retain complete declarations, with Unicode symbols
typeset mathematically and long lines wrapped. They are extracts from their
original namespaces and imported contexts, rather than separate standalone Lean
files. Unless stated otherwise, cyclic modules use an implicit natural number
\texttt{L} and the instance \texttt{[NeZero L]}. This excludes length zero, not
even or composite lengths. The generic cover modules instead use finite abelian
groups and the instances stated with their theorem signatures. The archive's
\nolinkurl{appendix-listings.json} binds each extract to its source hash and exact
line ranges. No mathematical proof source was changed to produce these listings.

\subsection{How the branches fit together}
The lower and upper bounds are proved independently and combined only after
their actual source-coordinate translations. The following diagram summarizes
the dependency structure; the graph and transpose branches establish the
additional clauses of Theorem~\ref{thm:main}.
\[
\begin{array}{ccc}
\text{Kernel and augmentation}&&\text{Base certificate and tree bound}\\
\downarrow&&\downarrow\\
\text{Full-span exclusion; weight }\leq11&&\text{Cover projection; weight }\geq12\\
\searrow&&\swarrow\\[-2pt]
\multicolumn{3}{c}{\text{Actual source words, spans and minima}}\\
\multicolumn{3}{c}{\downarrow}\\
\multicolumn{3}{c}{d_Q\leq11<12\leq d_C}
\end{array}
\]
\begin{center}\small
\begin{tabular}{@{}P{.30\textwidth}cP{.57\textwidth}@{}}
\toprule Branch&&Checked dependency chain\\\midrule
Quantum upper bound&$\longrightarrow$&
Coefficient kernel $\to$ augmentation separator $\to$ full-span nonmembership
$\to$ physical support at most eleven.\\[4pt]
Classical lower bound&$\longrightarrow$&
Base row reconstruction and finite weight certificate; Tanner-tree stopping bound
$\to$ parity projection $\to$ every nonzero upstairs word has weight at least twelve.\\[4pt]
Source distances&$\longrightarrow$&
Actual matrix actions and full spans $\to$ both CSS sectors and all-word weight
predicates $\to$ equality of the formal minima $\to$ $d_Q<d_C$.\\[4pt]
Graph and transpose properties&$\longrightarrow$&
Voltage exclusions plus explicit cycle edges; cyclic reachability;
transpose period identity $\to$ constant fibers $\to$ exact transpose weight.\\
\bottomrule
\end{tabular}
\end{center}
The trust boundary is the ordinary Lean kernel and the recorded standard axioms
\texttt{propext}, \texttt{Classical.choice} and \texttt{Quot.sound}. The finite
classical certificate is checked by Lean; external C++ and SAT exclusions are
absent from these branches. The separately audited generic follow-up lemmas are
not used to claim the unproved truncated-distance condition.

\subsection{Actual matrices and physical coordinates}
\nolinkurl{GeneralLP} gives the coefficient matrices for an arbitrary monomial
$3\times4$ exponent array. Its \texttt{Qubit} type is the disjoint union of the
$4\times4$ and $3\times3$ product blocks, each with one cyclic coordinate.
The \texttt{P} coefficient uses exactly $t=s+e$. Thus the definitions implement
\eqref{eq:source}, including the star's negative exponent, without a second lift.
The syndrome and two stabilizer actions displayed here expose every sign.
\leanlisting{source-objects}

\subsection{The complete row span, not just the generators}
The following equivalence quantifies over an arbitrary binary coefficient for
every lifted check. Its forward direction evaluates a general row combination
in each physical block; its reverse direction reconstructs that combination.
The \texttt{inSpanZ\_iff} theorem in the same file proves the corresponding
other-sector statement by transposing both square blocks. This is the bridge
needed before an augmentation obstruction can certify a logical class.
\leanlisting{source-span}

\subsection{The all-word distance predicate}
Here \texttt{Kernel} and \texttt{InSpan} quantify over actual source binary
coordinates, while \texttt{QuantumHasWeight} includes both logical sectors.
These are the predicates behind the minima, rather than a restricted search
through displayed words. Existence and the minimum equivalences are used in
Appendix~\ref{app:lean-assembly}.
\leanlisting{source-minimum}

\clearpage
\section{The logical obstruction: kernel, full-span separation and weight}\label{app:lean-obstruction}
This branch proves Theorem~\ref{thm:template}. The coefficient word is defined
for arbitrary parameters in \texttt{ZMod L}. In particular, monomial terms may
coincide, and the weight argument must allow characteristic-two cancellation.

\subsection{The word and its full coefficient kernel proof}
The matrix literal below is the coefficient version of \eqref{eq:UV}.
\texttt{delta a s} is one exactly when $s=a$. The proof checks all twelve
check types with arbitrary fiber coordinate; it does not sample lift sizes.
\leanlisting{parametric-word}

\subsection{Why the separator annihilates every stabilizer}
The augmentation is the two fiber sums of \eqref{eq:ell}. Translation permutes
a finite cyclic fiber, so both terms of a stabilizer give the same sum and
cancel. The second theorem extracts a left-block stabilizer representation
from membership in the entire physical row span. Together these statements
make the separator apply to arbitrary source stabilizer combinations.
\leanlisting{separator-span}

\subsection{Nontriviality of the actual source logical}
The source check action transports the coefficient identity to the matrix
kernel. If the word belonged to the stabilizer span, its left block would
have augmentation zero, contradicting the computed value one. The following
is the complete proof of that nonmembership, with no finite span search.
\leanlisting{logical-separation}

\subsection{Physical Hamming weight, including cancelling parameters}
The support list consists of eleven physical qubit coordinates. The first
proof shows the word vanishes outside that list; the second compares finite
support cardinalities. It deliberately uses list length as an upper bound,
so duplicate coordinates and cancellation cannot invalidate the theorem.
\leanlisting{physical-weight}

\subsection{The source-matched universal conclusion}
This theorem assembles actual CSS matrix commutation, full logical
nontriviality and physical coefficient weight. It is the same universal
parameter statement used in Theorem~\ref{thm:template}.
\leanlisting{source-obstruction}

\clearpage
\section{The classical lower bound: projection and the finite base}\label{app:lean-classical}
This branch separates ordinary binary parity projection from the image of
occupied support. The distinction is necessary when all projected parities
cancel, and avoids an odd-cover assumption.

\subsection{Cancellation gives two occupied coordinates per fiber}
The push-forward is a sum in \texttt{ZMod 2}. If it is zero, a nonempty
fiber support cannot have cardinality one. Summing the resulting two-per-fiber
bound gives twice the size of the occupied image. The full proof below
works for arbitrary finite source and target types, not just cyclic groups.
\leanlisting{cancelled-fibers}

\subsection{The occupied image is a nonempty stopping set}
The projected support theorem uses the actual monomial incidence. At a
touched lifted check, binary parity gives a second selected neighbor.
Its base column type differs from the first, so the neighbors remain
distinct after projection. The helper \texttt{stopping\_of\_partners}
turns these partners into the finite-set stopping predicate.
\leanlisting{stopping-projection}

\subsection{The complete cover-bound theorem}
Both alternatives now appear in a single full theorem and proof. When
the parity projection is nonzero, its base code distance supplies $d$;
when it is zero, the occupied image supplies $2s$. The printed hypotheses
make clear which base facts are required.
\leanlisting{cover-bound}

\subsection{What the eleven-fold certificate actually checks}
\nolinkurl{FamilyEnum11.encode} is an explicit XOR of thirteen $44$-bit
vectors. \nolinkurl{FamilySource11} proves
\texttt{kernel\_reconstruct}: every actual source-kernel vector equals
that encoding at its thirteen extracted free coordinates. It is an
identity derived from the source rows, not an assumed generator-span
description. \nolinkurl{FamilyEnum11.goodBlock} checks 128 consecutive
inputs, and \nolinkurl{FamilyDistance11.all_blocks} checks all 64 blocks.
Their complete source files, including the bit constants and row identities,
are linked in Appendix~\ref{app:lean-links}.

The following proofs turn the checked blocks into a theorem for an arbitrary
thirteen-bit input and then for every source-kernel vector. This displays
the certificate's soundness step rather than an external enumerator result.
\leanlisting{finite-certificate}

The next bridge packs actual coefficient functions into the bit-vector
certificate and identifies population count with physical Hamming weight.
The final theorem converts each source check equation into a certified row
dot product before applying the lower bound.
\leanlisting{finite-weight-bridge}

\subsection{The structural stopping bound}
The classical Tanner-tree argument is credited to Kelley et al.~\cite{Kelley}
in the main text. Here is its complete checked degree-three, no-four/no-six
implementation. It constructs six distinct selected variables explicitly;
the no-cycle hypotheses rule out each potential collision.
\leanlisting{tree-stopping}

\subsection{Application to every positive multiple of eleven}
Reduction modulo eleven supplies the additive map. The two premises of
the generic theorem are discharged by the actual base kernel certificate
and the independently proved structural stopping bound. This is the full
family application, including the exponent-projection equality.
\leanlisting{cover-application}

\clearpage
\section{Upper bounds, transpose distance and actual graph witnesses}\label{app:lean-geometry}

\subsection{A nonzero classical word of weight at most twenty}
The coefficient definition below is \eqref{eq:twenty}. Kernel membership
is proved for every positive $L$; nonzero membership uses reduction modulo
eleven. The final step applies the physical support bound to the actual
classical minimum. The helper \texttt{weight\_le} and its support list
are in the same linked source file.
\leanlisting{classical-upper}

\subsection{The transpose period identity}
The transpose boundary has the plus exponent $t+E_{ij}$. The following
proof writes down the precise combination of shifted check equations in
\eqref{eq:bezout}. The identity forces the third component to have period
one without semisimplicity, oddness or a rank-preservation hypothesis.
\leanlisting{transpose-period}

The intervening lemmas in \nolinkurl{TransposeConstant.lean} propagate
period one to constancy of all three fibers. The full equivalence below
includes the converse and the even-sum condition.
\leanlisting{transpose-constant}

\subsection{Exact weight of every nonzero transpose word}
The finite binary-triple fact \texttt{three\_even\_count} says that a
nonzero even triple has exactly two ones. The proof below then counts
the Cartesian product of those two occupied component types with all
$L$ fibers. The final conjunction includes existence, so the distance
statement is not vacuous.
\leanlisting{transpose-distance}

\subsection{Girth exactly eight}
The finite voltage tests and their reduction proofs exclude four- and
six-cycles. To complete exact girth, one also needs a simple eight-cycle
in the actual incidence. The following declarations give distinct
vertices and all eight edges separately. The main text explains their
assembly into the displayed cycle.
\leanlisting{girth-eight}

\subsection{Connectedness through actual Tanner walks}
The normalized zero row connects column types within a fiber; the unit
exponent supplies the following shift walk. The final proof combines
variable-to-variable walks and the zero-column incidences to cover all
variable/check endpoints. Both statements hold for every positive $L$.
\leanlisting{connected}

\clearpage
\section{Parameter transfer and the final source-distance statement}\label{app:lean-assembly}

\subsection{The arithmetic counterfamily}
This complete implication displays every premise of the general parameter
extension. An additive projection sends the parameters to $(1,5)$; the
conclusion includes the voltage exclusions, full partition exclusions,
actual classical-kernel lower bound and source logical upper witness.
The cyclic-congruence specialization is in the same source file.
\leanlisting{arithmetic-family}

\subsection{From block coordinates to every physical binary word}
The disjoint-union coordinate map is surjective, so every physical word
has two blocks. The first two equivalences transport logical-sector
membership and exact support weight. The next two proofs identify the
minima in both directions, using the corresponding weight equivalences.
This is the bridge that prevents an accidentally restricted minimum.
\leanlisting{minimum-bridge}

\subsection{The final distance drop}
All substantive ingredients now occur as checked dependencies. The
last theorem rewrites the actual source minima, applies the quantum
witness and classical lower theorem, and concludes strict separation.
The other clauses of Theorem~\ref{thm:main} are supplied by the graph,
transpose and upper-word proofs above; they are not extra assumptions
of this distance statement.
\leanlisting{final-distance}

\paragraph{Reproducing the checks.}
The fifty original modules and their full import closure were rebuilt
against the pinned cached dependencies as described in
Section~\ref{sec:formal}. The manuscript declaration audit checks every
named declaration linked here, including those added to the exposition.
The listing generator verifies that each source equals its published
Git blob, and that each printed declaration is a complete source range.
The archive contains that generator, the range/hash manifest, the
declaration audit and the portable rebuild instructions. The exact
typesetting extracts add no axioms or proof modules.

\clearpage\section{Complete current and immutable Lean URLs}\label{app:lean-links}
The following index includes every admitted proof, certificate and research import.
All immutable links use commit \texttt{98eb91cd9049e6a88887a5ba79150ad96a4a3e53}.
The declaration links in the text and preceding appendices go directly to source line ranges.
Current links may evolve; the immutable sources match the audited hashes.
The source archive and companion proof guide give the complete SHA-256 ledger and logs.
\par\smallskip\noindent\begin{minipage}{\linewidth}\small\setlength{\baselineskip}{11pt}
\textbf{Source file:} \nolinkurl{BinaryFinite.lean}\par\nopagebreak
Current: \url{https://github.com/alectionary/vibe-math/blob/master/formalization/distance-preservation/BinaryFinite.lean}\\
Immutable: \url{https://github.com/alectionary/vibe-math/blob/98eb91cd9049e6a88887a5ba79150ad96a4a3e53/formalization/distance-preservation/BinaryFinite.lean}\par\end{minipage}\par
\par\smallskip\noindent\begin{minipage}{\linewidth}\small\setlength{\baselineskip}{11pt}
\textbf{Source file:} \nolinkurl{BinaryWords.lean}\par\nopagebreak
Current: \url{https://github.com/alectionary/vibe-math/blob/master/formalization/distance-preservation/BinaryWords.lean}\\
Immutable: \url{https://github.com/alectionary/vibe-math/blob/98eb91cd9049e6a88887a5ba79150ad96a4a3e53/formalization/distance-preservation/BinaryWords.lean}\par\end{minipage}\par
\par\smallskip\noindent\begin{minipage}{\linewidth}\small\setlength{\baselineskip}{11pt}
\textbf{Source file:} \nolinkurl{CSS11.lean}\par\nopagebreak
Current: \url{https://github.com/alectionary/vibe-math/blob/master/formalization/distance-preservation/CSS11.lean}\\
Immutable: \url{https://github.com/alectionary/vibe-math/blob/98eb91cd9049e6a88887a5ba79150ad96a4a3e53/formalization/distance-preservation/CSS11.lean}\par\end{minipage}\par
\par\smallskip\noindent\begin{minipage}{\linewidth}\small\setlength{\baselineskip}{11pt}
\textbf{Source file:} \nolinkurl{ClassicalDistance11.lean}\par\nopagebreak
Current: \url{https://github.com/alectionary/vibe-math/blob/master/formalization/distance-preservation/ClassicalDistance11.lean}\\
Immutable: \url{https://github.com/alectionary/vibe-math/blob/98eb91cd9049e6a88887a5ba79150ad96a4a3e53/formalization/distance-preservation/ClassicalDistance11.lean}\par\end{minipage}\par
\par\smallskip\noindent\begin{minipage}{\linewidth}\small\setlength{\baselineskip}{11pt}
\textbf{Source file:} \nolinkurl{ClassicalEnum11.lean}\par\nopagebreak
Current: \url{https://github.com/alectionary/vibe-math/blob/master/formalization/distance-preservation/ClassicalEnum11.lean}\\
Immutable: \url{https://github.com/alectionary/vibe-math/blob/98eb91cd9049e6a88887a5ba79150ad96a4a3e53/formalization/distance-preservation/ClassicalEnum11.lean}\par\end{minipage}\par
\par\smallskip\noindent\begin{minipage}{\linewidth}\small\setlength{\baselineskip}{11pt}
\textbf{Source file:} \nolinkurl{ClassicalUpper20.lean}\par\nopagebreak
Current: \url{https://github.com/alectionary/vibe-math/blob/master/formalization/distance-preservation/ClassicalUpper20.lean}\\
Immutable: \url{https://github.com/alectionary/vibe-math/blob/98eb91cd9049e6a88887a5ba79150ad96a4a3e53/formalization/distance-preservation/ClassicalUpper20.lean}\par\end{minipage}\par
\par\smallskip\noindent\begin{minipage}{\linewidth}\small\setlength{\baselineskip}{11pt}
\textbf{Source file:} \nolinkurl{Counterexample11.lean}\par\nopagebreak
Current: \url{https://github.com/alectionary/vibe-math/blob/master/formalization/distance-preservation/Counterexample11.lean}\\
Immutable: \url{https://github.com/alectionary/vibe-math/blob/98eb91cd9049e6a88887a5ba79150ad96a4a3e53/formalization/distance-preservation/Counterexample11.lean}\par\end{minipage}\par
\par\smallskip\noindent\begin{minipage}{\linewidth}\small\setlength{\baselineskip}{11pt}
\textbf{Source file:} \nolinkurl{CoverDistance.lean}\par\nopagebreak
Current: \url{https://github.com/alectionary/vibe-math/blob/master/formalization/distance-preservation/CoverDistance.lean}\\
Immutable: \url{https://github.com/alectionary/vibe-math/blob/98eb91cd9049e6a88887a5ba79150ad96a4a3e53/formalization/distance-preservation/CoverDistance.lean}\par\end{minipage}\par
\par\smallskip\noindent\begin{minipage}{\linewidth}\small\setlength{\baselineskip}{11pt}
\textbf{Source file:} \nolinkurl{CoverWeight.lean}\par\nopagebreak
Current: \url{https://github.com/alectionary/vibe-math/blob/master/formalization/distance-preservation/CoverWeight.lean}\\
Immutable: \url{https://github.com/alectionary/vibe-math/blob/98eb91cd9049e6a88887a5ba79150ad96a4a3e53/formalization/distance-preservation/CoverWeight.lean}\par\end{minipage}\par
\par\smallskip\noindent\begin{minipage}{\linewidth}\small\setlength{\baselineskip}{11pt}
\textbf{Source file:} \nolinkurl{CyclicCover.lean}\par\nopagebreak
Current: \url{https://github.com/alectionary/vibe-math/blob/master/formalization/distance-preservation/CyclicCover.lean}\\
Immutable: \url{https://github.com/alectionary/vibe-math/blob/98eb91cd9049e6a88887a5ba79150ad96a4a3e53/formalization/distance-preservation/CyclicCover.lean}\par\end{minipage}\par
\par\smallskip\noindent\begin{minipage}{\linewidth}\small\setlength{\baselineskip}{11pt}
\textbf{Source file:} \nolinkurl{FamilyBaseDistance.lean}\par\nopagebreak
Current: \url{https://github.com/alectionary/vibe-math/blob/master/formalization/distance-preservation/FamilyBaseDistance.lean}\\
Immutable: \url{https://github.com/alectionary/vibe-math/blob/98eb91cd9049e6a88887a5ba79150ad96a4a3e53/formalization/distance-preservation/FamilyBaseDistance.lean}\par\end{minipage}\par
\par\smallskip\noindent\begin{minipage}{\linewidth}\small\setlength{\baselineskip}{11pt}
\textbf{Source file:} \nolinkurl{FamilyBaseGraph.lean}\par\nopagebreak
Current: \url{https://github.com/alectionary/vibe-math/blob/master/formalization/distance-preservation/FamilyBaseGraph.lean}\\
Immutable: \url{https://github.com/alectionary/vibe-math/blob/98eb91cd9049e6a88887a5ba79150ad96a4a3e53/formalization/distance-preservation/FamilyBaseGraph.lean}\par\end{minipage}\par
\par\smallskip\noindent\begin{minipage}{\linewidth}\small\setlength{\baselineskip}{11pt}
\textbf{Source file:} \nolinkurl{FamilyBounds.lean}\par\nopagebreak
Current: \url{https://github.com/alectionary/vibe-math/blob/master/formalization/distance-preservation/FamilyBounds.lean}\\
Immutable: \url{https://github.com/alectionary/vibe-math/blob/98eb91cd9049e6a88887a5ba79150ad96a4a3e53/formalization/distance-preservation/FamilyBounds.lean}\par\end{minipage}\par
\par\smallskip\noindent\begin{minipage}{\linewidth}\small\setlength{\baselineskip}{11pt}
\textbf{Source file:} \nolinkurl{FamilyConnected.lean}\par\nopagebreak
Current: \url{https://github.com/alectionary/vibe-math/blob/master/formalization/distance-preservation/FamilyConnected.lean}\\
Immutable: \url{https://github.com/alectionary/vibe-math/blob/98eb91cd9049e6a88887a5ba79150ad96a4a3e53/formalization/distance-preservation/FamilyConnected.lean}\par\end{minipage}\par
\par\smallskip\noindent\begin{minipage}{\linewidth}\small\setlength{\baselineskip}{11pt}
\textbf{Source file:} \nolinkurl{FamilyCoverAlternative.lean}\par\nopagebreak
Current: \url{https://github.com/alectionary/vibe-math/blob/master/formalization/distance-preservation/FamilyCoverAlternative.lean}\\
Immutable: \url{https://github.com/alectionary/vibe-math/blob/98eb91cd9049e6a88887a5ba79150ad96a4a3e53/formalization/distance-preservation/FamilyCoverAlternative.lean}\par\end{minipage}\par
\par\smallskip\noindent\begin{minipage}{\linewidth}\small\setlength{\baselineskip}{11pt}
\textbf{Source file:} \nolinkurl{FamilyDistance11.lean}\par\nopagebreak
Current: \url{https://github.com/alectionary/vibe-math/blob/master/formalization/distance-preservation/FamilyDistance11.lean}\\
Immutable: \url{https://github.com/alectionary/vibe-math/blob/98eb91cd9049e6a88887a5ba79150ad96a4a3e53/formalization/distance-preservation/FamilyDistance11.lean}\par\end{minipage}\par
\par\smallskip\noindent\begin{minipage}{\linewidth}\small\setlength{\baselineskip}{11pt}
\textbf{Source file:} \nolinkurl{FamilyDotRows11.lean}\par\nopagebreak
Current: \url{https://github.com/alectionary/vibe-math/blob/master/formalization/distance-preservation/FamilyDotRows11.lean}\\
Immutable: \url{https://github.com/alectionary/vibe-math/blob/98eb91cd9049e6a88887a5ba79150ad96a4a3e53/formalization/distance-preservation/FamilyDotRows11.lean}\par\end{minipage}\par
\par\smallskip\noindent\begin{minipage}{\linewidth}\small\setlength{\baselineskip}{11pt}
\textbf{Source file:} \nolinkurl{FamilyEnum11.lean}\par\nopagebreak
Current: \url{https://github.com/alectionary/vibe-math/blob/master/formalization/distance-preservation/FamilyEnum11.lean}\\
Immutable: \url{https://github.com/alectionary/vibe-math/blob/98eb91cd9049e6a88887a5ba79150ad96a4a3e53/formalization/distance-preservation/FamilyEnum11.lean}\par\end{minipage}\par
\par\smallskip\noindent\begin{minipage}{\linewidth}\small\setlength{\baselineskip}{11pt}
\textbf{Source file:} \nolinkurl{FamilyGirth.lean}\par\nopagebreak
Current: \url{https://github.com/alectionary/vibe-math/blob/master/formalization/distance-preservation/FamilyGirth.lean}\\
Immutable: \url{https://github.com/alectionary/vibe-math/blob/98eb91cd9049e6a88887a5ba79150ad96a4a3e53/formalization/distance-preservation/FamilyGirth.lean}\par\end{minipage}\par
\par\smallskip\noindent\begin{minipage}{\linewidth}\small\setlength{\baselineskip}{11pt}
\textbf{Source file:} \nolinkurl{FamilyNoShort.lean}\par\nopagebreak
Current: \url{https://github.com/alectionary/vibe-math/blob/master/formalization/distance-preservation/FamilyNoShort.lean}\\
Immutable: \url{https://github.com/alectionary/vibe-math/blob/98eb91cd9049e6a88887a5ba79150ad96a4a3e53/formalization/distance-preservation/FamilyNoShort.lean}\par\end{minipage}\par
\par\smallskip\noindent\begin{minipage}{\linewidth}\small\setlength{\baselineskip}{11pt}
\textbf{Source file:} \nolinkurl{FamilyRing.lean}\par\nopagebreak
Current: \url{https://github.com/alectionary/vibe-math/blob/master/formalization/distance-preservation/FamilyRing.lean}\\
Immutable: \url{https://github.com/alectionary/vibe-math/blob/98eb91cd9049e6a88887a5ba79150ad96a4a3e53/formalization/distance-preservation/FamilyRing.lean}\par\end{minipage}\par
\par\smallskip\noindent\begin{minipage}{\linewidth}\small\setlength{\baselineskip}{11pt}
\textbf{Source file:} \nolinkurl{FamilySource11.lean}\par\nopagebreak
Current: \url{https://github.com/alectionary/vibe-math/blob/master/formalization/distance-preservation/FamilySource11.lean}\\
Immutable: \url{https://github.com/alectionary/vibe-math/blob/98eb91cd9049e6a88887a5ba79150ad96a4a3e53/formalization/distance-preservation/FamilySource11.lean}\par\end{minipage}\par
\par\smallskip\noindent\begin{minipage}{\linewidth}\small\setlength{\baselineskip}{11pt}
\textbf{Source file:} \nolinkurl{FamilyStopping11.lean}\par\nopagebreak
Current: \url{https://github.com/alectionary/vibe-math/blob/master/formalization/distance-preservation/FamilyStopping11.lean}\\
Immutable: \url{https://github.com/alectionary/vibe-math/blob/98eb91cd9049e6a88887a5ba79150ad96a4a3e53/formalization/distance-preservation/FamilyStopping11.lean}\par\end{minipage}\par
\par\smallskip\noindent\begin{minipage}{\linewidth}\small\setlength{\baselineskip}{11pt}
\textbf{Source file:} \nolinkurl{FamilyTheorem.lean}\par\nopagebreak
Current: \url{https://github.com/alectionary/vibe-math/blob/master/formalization/distance-preservation/FamilyTheorem.lean}\\
Immutable: \url{https://github.com/alectionary/vibe-math/blob/98eb91cd9049e6a88887a5ba79150ad96a4a3e53/formalization/distance-preservation/FamilyTheorem.lean}\par\end{minipage}\par
\par\smallskip\noindent\begin{minipage}{\linewidth}\small\setlength{\baselineskip}{11pt}
\textbf{Source file:} \nolinkurl{FamilyTranspose.lean}\par\nopagebreak
Current: \url{https://github.com/alectionary/vibe-math/blob/master/formalization/distance-preservation/FamilyTranspose.lean}\\
Immutable: \url{https://github.com/alectionary/vibe-math/blob/98eb91cd9049e6a88887a5ba79150ad96a4a3e53/formalization/distance-preservation/FamilyTranspose.lean}\par\end{minipage}\par
\par\smallskip\noindent\begin{minipage}{\linewidth}\small\setlength{\baselineskip}{11pt}
\textbf{Source file:} \nolinkurl{FamilyTreeStopping.lean}\par\nopagebreak
Current: \url{https://github.com/alectionary/vibe-math/blob/master/formalization/distance-preservation/FamilyTreeStopping.lean}\\
Immutable: \url{https://github.com/alectionary/vibe-math/blob/98eb91cd9049e6a88887a5ba79150ad96a4a3e53/formalization/distance-preservation/FamilyTreeStopping.lean}\par\end{minipage}\par
\par\smallskip\noindent\begin{minipage}{\linewidth}\small\setlength{\baselineskip}{11pt}
\textbf{Source file:} \nolinkurl{FamilyWords.lean}\par\nopagebreak
Current: \url{https://github.com/alectionary/vibe-math/blob/master/formalization/distance-preservation/FamilyWords.lean}\\
Immutable: \url{https://github.com/alectionary/vibe-math/blob/98eb91cd9049e6a88887a5ba79150ad96a4a3e53/formalization/distance-preservation/FamilyWords.lean}\par\end{minipage}\par
\par\smallskip\noindent\begin{minipage}{\linewidth}\small\setlength{\baselineskip}{11pt}
\textbf{Source file:} \nolinkurl{FinalAudit.lean}\par\nopagebreak
Current: \url{https://github.com/alectionary/vibe-math/blob/master/formalization/distance-preservation/FinalAudit.lean}\\
Immutable: \url{https://github.com/alectionary/vibe-math/blob/98eb91cd9049e6a88887a5ba79150ad96a4a3e53/formalization/distance-preservation/FinalAudit.lean}\par\end{minipage}\par
\par\smallskip\noindent\begin{minipage}{\linewidth}\small\setlength{\baselineskip}{11pt}
\textbf{Source file:} \nolinkurl{GeneralLP.lean}\par\nopagebreak
Current: \url{https://github.com/alectionary/vibe-math/blob/master/formalization/distance-preservation/GeneralLP.lean}\\
Immutable: \url{https://github.com/alectionary/vibe-math/blob/98eb91cd9049e6a88887a5ba79150ad96a4a3e53/formalization/distance-preservation/GeneralLP.lean}\par\end{minipage}\par
\par\smallskip\noindent\begin{minipage}{\linewidth}\small\setlength{\baselineskip}{11pt}
\textbf{Source file:} \nolinkurl{GeneralLPPhysical.lean}\par\nopagebreak
Current: \url{https://github.com/alectionary/vibe-math/blob/master/formalization/distance-preservation/GeneralLPPhysical.lean}\\
Immutable: \url{https://github.com/alectionary/vibe-math/blob/98eb91cd9049e6a88887a5ba79150ad96a4a3e53/formalization/distance-preservation/GeneralLPPhysical.lean}\par\end{minipage}\par
\par\smallskip\noindent\begin{minipage}{\linewidth}\small\setlength{\baselineskip}{11pt}
\textbf{Source file:} \nolinkurl{Hypotheses11.lean}\par\nopagebreak
Current: \url{https://github.com/alectionary/vibe-math/blob/master/formalization/distance-preservation/Hypotheses11.lean}\\
Immutable: \url{https://github.com/alectionary/vibe-math/blob/98eb91cd9049e6a88887a5ba79150ad96a4a3e53/formalization/distance-preservation/Hypotheses11.lean}\par\end{minipage}\par
\par\smallskip\noindent\begin{minipage}{\linewidth}\small\setlength{\baselineskip}{11pt}
\textbf{Source file:} \nolinkurl{ParametricProjection.lean}\par\nopagebreak
Current: \url{https://github.com/alectionary/vibe-math/blob/master/formalization/distance-preservation/ParametricProjection.lean}\\
Immutable: \url{https://github.com/alectionary/vibe-math/blob/98eb91cd9049e6a88887a5ba79150ad96a4a3e53/formalization/distance-preservation/ParametricProjection.lean}\par\end{minipage}\par
\par\smallskip\noindent\begin{minipage}{\linewidth}\small\setlength{\baselineskip}{11pt}
\textbf{Source file:} \nolinkurl{ParametricRing.lean}\par\nopagebreak
Current: \url{https://github.com/alectionary/vibe-math/blob/master/formalization/distance-preservation/ParametricRing.lean}\\
Immutable: \url{https://github.com/alectionary/vibe-math/blob/98eb91cd9049e6a88887a5ba79150ad96a4a3e53/formalization/distance-preservation/ParametricRing.lean}\par\end{minipage}\par
\par\smallskip\noindent\begin{minipage}{\linewidth}\small\setlength{\baselineskip}{11pt}
\textbf{Source file:} \nolinkurl{ParametricSource.lean}\par\nopagebreak
Current: \url{https://github.com/alectionary/vibe-math/blob/master/formalization/distance-preservation/ParametricSource.lean}\\
Immutable: \url{https://github.com/alectionary/vibe-math/blob/98eb91cd9049e6a88887a5ba79150ad96a4a3e53/formalization/distance-preservation/ParametricSource.lean}\par\end{minipage}\par
\par\smallskip\noindent\begin{minipage}{\linewidth}\small\setlength{\baselineskip}{11pt}
\textbf{Source file:} \nolinkurl{ParametricSpecialization.lean}\par\nopagebreak
Current: \url{https://github.com/alectionary/vibe-math/blob/master/formalization/distance-preservation/ParametricSpecialization.lean}\\
Immutable: \url{https://github.com/alectionary/vibe-math/blob/98eb91cd9049e6a88887a5ba79150ad96a4a3e53/formalization/distance-preservation/ParametricSpecialization.lean}\par\end{minipage}\par
\par\smallskip\noindent\begin{minipage}{\linewidth}\small\setlength{\baselineskip}{11pt}
\textbf{Source file:} \nolinkurl{ParametricWords.lean}\par\nopagebreak
Current: \url{https://github.com/alectionary/vibe-math/blob/master/formalization/distance-preservation/ParametricWords.lean}\\
Immutable: \url{https://github.com/alectionary/vibe-math/blob/98eb91cd9049e6a88887a5ba79150ad96a4a3e53/formalization/distance-preservation/ParametricWords.lean}\par\end{minipage}\par
\par\smallskip\noindent\begin{minipage}{\linewidth}\small\setlength{\baselineskip}{11pt}
\textbf{Source file:} \nolinkurl{QuantumRows11.lean}\par\nopagebreak
Current: \url{https://github.com/alectionary/vibe-math/blob/master/formalization/distance-preservation/QuantumRows11.lean}\\
Immutable: \url{https://github.com/alectionary/vibe-math/blob/98eb91cd9049e6a88887a5ba79150ad96a4a3e53/formalization/distance-preservation/QuantumRows11.lean}\par\end{minipage}\par
\par\smallskip\noindent\begin{minipage}{\linewidth}\small\setlength{\baselineskip}{11pt}
\textbf{Source file:} \nolinkurl{QuantumWitness11.lean}\par\nopagebreak
Current: \url{https://github.com/alectionary/vibe-math/blob/master/formalization/distance-preservation/QuantumWitness11.lean}\\
Immutable: \url{https://github.com/alectionary/vibe-math/blob/98eb91cd9049e6a88887a5ba79150ad96a4a3e53/formalization/distance-preservation/QuantumWitness11.lean}\par\end{minipage}\par
\par\smallskip\noindent\begin{minipage}{\linewidth}\small\setlength{\baselineskip}{11pt}
\textbf{Source file:} \nolinkurl{SmallStopping.lean}\par\nopagebreak
Current: \url{https://github.com/alectionary/vibe-math/blob/master/formalization/distance-preservation/SmallStopping.lean}\\
Immutable: \url{https://github.com/alectionary/vibe-math/blob/98eb91cd9049e6a88887a5ba79150ad96a4a3e53/formalization/distance-preservation/SmallStopping.lean}\par\end{minipage}\par
\par\smallskip\noindent\begin{minipage}{\linewidth}\small\setlength{\baselineskip}{11pt}
\textbf{Source file:} \nolinkurl{Source11.lean}\par\nopagebreak
Current: \url{https://github.com/alectionary/vibe-math/blob/master/formalization/distance-preservation/Source11.lean}\\
Immutable: \url{https://github.com/alectionary/vibe-math/blob/98eb91cd9049e6a88887a5ba79150ad96a4a3e53/formalization/distance-preservation/Source11.lean}\par\end{minipage}\par
\par\smallskip\noindent\begin{minipage}{\linewidth}\small\setlength{\baselineskip}{11pt}
\textbf{Source file:} \nolinkurl{SourceCoordinates.lean}\par\nopagebreak
Current: \url{https://github.com/alectionary/vibe-math/blob/master/formalization/distance-preservation/SourceCoordinates.lean}\\
Immutable: \url{https://github.com/alectionary/vibe-math/blob/98eb91cd9049e6a88887a5ba79150ad96a4a3e53/formalization/distance-preservation/SourceCoordinates.lean}\par\end{minipage}\par
\par\smallskip\noindent\begin{minipage}{\linewidth}\small\setlength{\baselineskip}{11pt}
\textbf{Source file:} \nolinkurl{SourceDistance.lean}\par\nopagebreak
Current: \url{https://github.com/alectionary/vibe-math/blob/master/formalization/distance-preservation/SourceDistance.lean}\\
Immutable: \url{https://github.com/alectionary/vibe-math/blob/98eb91cd9049e6a88887a5ba79150ad96a4a3e53/formalization/distance-preservation/SourceDistance.lean}\par\end{minipage}\par
\par\smallskip\noindent\begin{minipage}{\linewidth}\small\setlength{\baselineskip}{11pt}
\textbf{Source file:} \nolinkurl{SourceModel.lean}\par\nopagebreak
Current: \url{https://github.com/alectionary/vibe-math/blob/master/formalization/distance-preservation/SourceModel.lean}\\
Immutable: \url{https://github.com/alectionary/vibe-math/blob/98eb91cd9049e6a88887a5ba79150ad96a4a3e53/formalization/distance-preservation/SourceModel.lean}\par\end{minipage}\par
\par\smallskip\noindent\begin{minipage}{\linewidth}\small\setlength{\baselineskip}{11pt}
\textbf{Source file:} \nolinkurl{SourceSectors.lean}\par\nopagebreak
Current: \url{https://github.com/alectionary/vibe-math/blob/master/formalization/distance-preservation/SourceSectors.lean}\\
Immutable: \url{https://github.com/alectionary/vibe-math/blob/98eb91cd9049e6a88887a5ba79150ad96a4a3e53/formalization/distance-preservation/SourceSectors.lean}\par\end{minipage}\par
\par\smallskip\noindent\begin{minipage}{\linewidth}\small\setlength{\baselineskip}{11pt}
\textbf{Source file:} \nolinkurl{SourceSymmetry.lean}\par\nopagebreak
Current: \url{https://github.com/alectionary/vibe-math/blob/master/formalization/distance-preservation/SourceSymmetry.lean}\\
Immutable: \url{https://github.com/alectionary/vibe-math/blob/98eb91cd9049e6a88887a5ba79150ad96a4a3e53/formalization/distance-preservation/SourceSymmetry.lean}\par\end{minipage}\par
\par\smallskip\noindent\begin{minipage}{\linewidth}\small\setlength{\baselineskip}{11pt}
\textbf{Source file:} \nolinkurl{SourceWordBridge.lean}\par\nopagebreak
Current: \url{https://github.com/alectionary/vibe-math/blob/master/formalization/distance-preservation/SourceWordBridge.lean}\\
Immutable: \url{https://github.com/alectionary/vibe-math/blob/98eb91cd9049e6a88887a5ba79150ad96a4a3e53/formalization/distance-preservation/SourceWordBridge.lean}\par\end{minipage}\par
\par\smallskip\noindent\begin{minipage}{\linewidth}\small\setlength{\baselineskip}{11pt}
\textbf{Source file:} \nolinkurl{TannerSix.lean}\par\nopagebreak
Current: \url{https://github.com/alectionary/vibe-math/blob/master/formalization/distance-preservation/TannerSix.lean}\\
Immutable: \url{https://github.com/alectionary/vibe-math/blob/98eb91cd9049e6a88887a5ba79150ad96a4a3e53/formalization/distance-preservation/TannerSix.lean}\par\end{minipage}\par
\par\smallskip\noindent\begin{minipage}{\linewidth}\small\setlength{\baselineskip}{11pt}
\textbf{Source file:} \nolinkurl{TransposeConstant.lean}\par\nopagebreak
Current: \url{https://github.com/alectionary/vibe-math/blob/master/formalization/distance-preservation/TransposeConstant.lean}\\
Immutable: \url{https://github.com/alectionary/vibe-math/blob/98eb91cd9049e6a88887a5ba79150ad96a4a3e53/formalization/distance-preservation/TransposeConstant.lean}\par\end{minipage}\par
\par\smallskip\noindent\begin{minipage}{\linewidth}\small\setlength{\baselineskip}{11pt}
\textbf{Source file:} \nolinkurl{TransposeDistance.lean}\par\nopagebreak
Current: \url{https://github.com/alectionary/vibe-math/blob/master/formalization/distance-preservation/TransposeDistance.lean}\\
Immutable: \url{https://github.com/alectionary/vibe-math/blob/98eb91cd9049e6a88887a5ba79150ad96a4a3e53/formalization/distance-preservation/TransposeDistance.lean}\par\end{minipage}\par
\par\smallskip\noindent\begin{minipage}{\linewidth}\small\setlength{\baselineskip}{11pt}
\textbf{Source file:} \nolinkurl{VoltageCycles.lean}\par\nopagebreak
Current: \url{https://github.com/alectionary/vibe-math/blob/master/formalization/distance-preservation/VoltageCycles.lean}\\
Immutable: \url{https://github.com/alectionary/vibe-math/blob/98eb91cd9049e6a88887a5ba79150ad96a4a3e53/formalization/distance-preservation/VoltageCycles.lean}\par\end{minipage}\par


\clearpage
\section{Reference dates and evidence boundaries}\label{app:references}
The chronology in Section~\ref{sec:answer} uses the manuscript date
25 October 2026 in UTC. Calendar subtraction from 10 March 2025 gives
594 days and from 10 June 2026 gives 137 days. Only year precision is
used for the inspected ISIT 2004 source; its age is therefore stated
approximately. The arXiv histories inspected for \cite{RDV,FHM} list
v1 only. These checks do not establish global unresolved status.
Both July preprints were inspected in the v1 copies recorded in the
evidence ledger, which carry an internal manuscript date of 24 August
2026. The corresponding deltas are 87 and 62 days. The current abstract
page for \cite{Bhardwaj} additionally lists v2 dated 2 October 2026;
that later version was identified during report preparation and is
retained as a version-comparison gap, not treated as a screened full text.
The comparison above is explicitly to the inspected v1 theorem.

The complete investigation ledger is
\href{../prior-art.json}{\texttt{prior-art.json}} and the semantic
translation record is \href{../SOURCE-TO-PROOF.txt}{\texttt{SOURCE-TO-PROOF.txt}}.
The prior-art audit includes primary papers, author references, public
GitHub code and signed-in Edge searches of GitHub and X. It records
queries, known-positive controls, caps and inspected results. Later
native-browser access failed when the Mac locked; indexed search
did not replace complete native coverage. Some discovered leads
were available only as snippets and were not used as theorem evidence.
The precise unresolved leads and access limitations remain in that
ledger. The searches support concrete source comparisons, not a
claim that no earlier result exists anywhere.

The supporting research record separates general Lean proofs, concrete
Lean certificates, exact finite computations and sampled diagnostics.
In particular, negative C++ searches are external computations. A
completion-shortcut bug was repaired, with an old-fails/new-passes
regression, 1,235 small full-word comparisons and replay of every
original saved input. Only corrected replay records support current
negative search statements. These computations are absent from the
theorem dependency closure. The separate SAT attempt returned
\texttt{UNKNOWN}, not a distance certificate. None of these
investigation details changes the proved counterfamily.

The companion \href{README.txt}{\texttt{README.txt}} gives artifact
layout and reproduction commands. The manuscript source archive
contains the fifty admitted Lean sources, their checked build records,
the dependency pins, the source-to-proof guide, evidence ledgers and
the manuscript. Content hashes identify the local review snapshot.
The investigation duration and subsequent manuscript preparation are
recorded there separately; they are not source chronology claims.

\begin{thebibliography}{9}
\bibitem{RDV}
N. Raveendran, D. Declercq and B. Vasi\'c,
\emph{On the Minimum Distances of Finite-Length Lifted Product Quantum
LDPC Codes}, arXiv:2503.07567v1, 10 March 2025.
\href{https://arxiv.org/html/2503.07567v1}{Full text};
Eq.~(9), Definitions 2--5, Theorem 2, Sections V-B and VI.

\bibitem{FHM}
R. Freij-Hollanti, K. D. Morris and P. \v Sapokait\.{e},
\emph{Graphical Analysis of Lifted Product Code Constructions},
arXiv:2606.11987v1, 10 June 2026.
\href{https://arxiv.org/html/2606.11987v1}{Full text};
sector graph isomorphism, connectivity and Section 5.

\bibitem{Kelley}
C. Kelley, D. Sridhara, J. Xu and J. Rosenthal,
\emph{Pseudocodeword Weights and Stopping Sets},
Proceedings of the IEEE International Symposium on Information Theory,
2004. Theorem III.1 and its following note; Sections II--III.
\href{https://user.math.uzh.ch/kelley/isit2004_Chicago.pdf}{Author-hosted paper}.

\bibitem{Bhardwaj}
A. Bhardwaj et al., \emph{High-rate qLDPC processors},
arXiv:2607.28795v1, first submitted 30 July 2026; inspected copy
internally dated 24 August 2026. Appendix G, Theorem 8 and Remark 12.
\href{https://arxiv.org/html/2607.28795v1}{Inspected version};
\href{https://arxiv.org/abs/2607.28795}{version history} (later v2 identified).

\bibitem{Hirasaki}
Y. Hirasaki and J. Y. Lee, \emph{Lifting Lifted Product Codes},
arXiv:2607.28621v1, first submitted 30 July 2026; inspected copy
internally dated 24 August 2026. Section II and Corollary 6.
\href{https://arxiv.org/html/2607.28621v1}{Full text}.
\end{thebibliography}
\end{document}
